\chapter{Boundary Values of Harmonic Functions}
\label{ch:c2}

A basic tool in the study of a function, $f$, on $E_n$ is obtained by the passage to a harmonic function\index{harmonic function@harmonic function}, defined on the $(n+1)$-dimensional upper half-space\index{upper half space@upper half space}, which has $f$ as its boundary value. In the first section we set down some of the useful basic facts concerning harmonic functions of several variables. The connection between an $f$, defined on $E_n$, and a harmonic function of $(n+1)$ variables (its Poisson integral) is studied in Section 2. Sections 3 and 4 deal with nontangential convergence at the boundary (for harmonic functions) and with harmonic majorization. While these notions are interesting in themselves, they will give us some important tools used in Chapters III and VI.

\section{Basic Properties of Harmonic Functions}

We assume the reader is familiar with the properties of harmonic functions that are usually presented in a course in complex variables. Many of these results extend to higher dimensions; however, we can no longer rely on the theory of analytic functions (as is often done in two dimensions) to establish these properties. In this section, therefore, we develop some of the basic theory of harmonic functions of n variables.

A function, $u$, defined in a domain (an open connected subset of $E_n$) is said to be harmonic if it is twice differentiable (i.e., its second partial derivatives exist and are continuous) and satisfies the Laplace\index{Laplace equation@Laplace equation} equation $\Delta u = \sum_{k=1}^{n} \partial^2 u / \partial x_k^2 = 0$ at every point of the domain. From the first two of the following examples of harmonic functions we can easily deduce that the Abel means and Poisson integrals introduced in the first section of Chapter I are harmonic functions:

(1) For $t \in E_n$ fixed $e^{-2\pi|t|y} e^{2\pi i x \cdot t} = u(x, y) = u(x_1, x_2, \ldots, x_n, y)$ is harmonic in $E_{n+1}$;

(2) $P(x, y) = c_n[y/(|x|^2 + y^2)^{(n+1)/2}]$ is also harmonic in $E_{n+1}^+$ = $\{(x, y) = (x_1, \ldots, x_n, y) \in E_{n+1}; y > 0\}$(this domain will often be called the upper half-space) \footnote{We denote the points of $E_{n+1}^+$ by the ordered pairs $(x, y)$, where $x \in E_n$ and $y$ is a positive real, in order to emphasize the fact that this domain is a generalization of the upper half-plane (the case $n = 1$).};

(3) If $n \geq 3$ then $u(x) = |x|^{2-n}$ is harmonic in any domain not containing the origin. The corresponding example in two dimensions is $u(x) = \log |x|$. In a sense that will become clearer later, these functions ``generate'' all harmonic functions;

(4) We shall show later (Theorem \ref{thm:c2-1-7}) that a harmonic function is infinitely differentiable, and, in fact, is real analytic (see \eqref{frp:c4-5-5} Chapter IV). It is clear that any partial derivative of a harmonic function is also harmonic. As an illustration we observe that $P(x, y)$ is the partial derivative with respect to $y$ of $c_n/(1 - n) \cdot 1/(|x|^2 + y^2)^{(n-1)/2}$, when $n > 1$, and of $c_n \log(|x|^2 + y^2)^{1/2}$, when $n$ is 1.

(5) If $u$ is harmonic in a domain $\mathcal{D}$ then $\tau_h u$ is harmonic in $\mathcal{D} + h = \{x + h; x \in \mathcal{D}\}$. That is, harmonicity is preserved under translations. Similarly, harmonicity is preserved under rotations and dilations. (This can be computed directly, but will be more evident after we establish the characterization of harmonic functions in terms of the mean value\index{mean value@mean value} property---see Theorem \ref{thm:c2-1-1} below);\index{mean value@mean value!theorem@theorem}

(6) $\sin \sqrt{n} y \prod_{k=1}^{n} \cosh x_k = u(x, y)$ is harmonic in $E_{n+1} = \{(x, y); x \in E_n, y \in E_1\}$.

We shall denote the mean value of a function $u$ taken over the surface of a sphere of radius $r > 0$ about the point $x$ by the symbols $\mathcal{M}(r) = \mathcal{M}_{x,u}(r)$ (depending on the necessity of showing its dependence on $u$, $x$ and $r$). That is, letting $\Sigma$ denote the surface of the unit sphere in $E_n$ and $dt'$ the element of surface area on $\Sigma$ (this was the notation used in the proof of Lemma \ref{lem:c1-1-17} in Chapter I)

\[
\mathcal{M}_{x,u}(r)=\frac{1}{\omega_{n-1}}\intop_{\Sigma}u(x+r t^{\prime}) d t^{\prime}.
\]

\begin{theorem}
\label{thm:c2-1-1}
(Mean-Value Theorem for Harmonic Functions). If $u$ is harmonic in a domain $\mathcal{D}$ and if the sphere of radius $r_0$ about $x\in\mathcal{D}$ is contained in $\mathcal{D}$ then, for $0<r\leq r_0$,
\[
u(x)=\mathcal{M}_{x,u}(r).
\]

\end{theorem}

\textsc{Proof.}\ Because of (5) we see that we can assume $x = 0$. Let $\mathcal{E}$ be a subdomain with a boundary $\partial\mathcal{E} \subset \mathcal{D}$ that is sufficiently smooth. Then, if we apply Green's theorem to the functions $u$ and 1 in $\mathcal{E}$, we obtain

\begin{equation}\label{eq:c2-1-2}
\int\limits_{\partial\mathcal{E}}\frac{\partial u}{\partial n}d s=0,
\end{equation}

where $\partial/\partial n$ is differentiation in the direction of the outward directed normal to $\partial\mathcal{E}$ and $ds$ is the element of surface area on $\partial\mathcal{E}$.

Let $\Sigma_\varepsilon$ and $\Sigma_r$, $0 < \varepsilon < r \leq r_0$, be the surfaces of the spheres of radii $\varepsilon$ and $r$ about 0. Applying Green's theorem to the functions $u$ and

\[
v(x)=\left\{\begin{aligned}&\left|x\right|^{-(n-2)}&\quad&\text{if } n>2\\ &\log\left|x\right|&\quad&\text{if } n=2\end{aligned}\right.
\]

in the shell domain between these two spheres we obtain \footnote{The proof will be written out for the case n > 2. Obvious modifications give us the result when n = 2. When n = 1 the theorem is obvious since the harmonic functions are the linear ones in this case.}

\[
0=\biggl(\int_{\Sigma_{r}}-\int_{\Sigma_{\varepsilon}}\biggr)(-(n-2)|x|^{-(n-1)}u) d s-\biggl\{\int_{\Sigma_{r}}v\frac{\partial u}{\partial n}d s+\int_{\Sigma_{\varepsilon}}v\frac{\partial u}{\partial n}d s\biggr\}.
\]

But $v$ is constant on $\Sigma_r$ (where it equals $r^{-(n-2)}$) and on $\Sigma_\varepsilon$ (where it equals $\varepsilon^{-(n-2)}$). Since $u$ is harmonic in $\mathcal{D} \supset \{x \in E_n; |x| \leq r_0\}$, applying \eqref{eq:c2-1-2} to $\partial \mathcal{E} = \Sigma_r$ and to $\partial \mathcal{E} = \Sigma_\varepsilon$, we see that the term in curly brackets must be 0. Consequently

\[
\varepsilon^{1-n}\int_{\Sigma_{\varepsilon}}^{\displaystyle\sum_{\varepsilon}}u d s = r^{1-n}\int_{\Sigma_{r}}^{\displaystyle\sum_{\varepsilon}}u d s.
\]

Hence,

\[
\begin{aligned}\mathcal{M}_{0,u}(r)&=\frac{1}{\omega_{n-1}}\int\limits_{\Sigma}u(rt^{\prime}) dt^{\prime}=\frac{1}{\omega_{n-1}r^{n-1}}\int\limits_{\Sigma_{r}}u ds\\&=\frac{1}{\omega_{n-1}\varepsilon^{n-1}}\int\limits_{\Sigma_{\varepsilon}}u ds.\end{aligned}
\]

But the last expression obviously tends to $u(0)$ as $\varepsilon \to 0$; this proves the theorem.

\begin{corollary}
\label{cor:c2-1-3}
(Maximum Principle\index{maximum principle@maximum principle} for Harmonic Functions). Suppose a real-valued harmonic function $u$, defined in a domain $\mathcal{D}$, satisfies $A = \sup_{x \in \mathcal{D}} u(x) < \infty$; then $u(x) < A$ for all $x \in \mathcal{D}$, provided $u$ is not a constant function.
\end{corollary}

\textsc{Proof.}\ Suppose $u(x) = A$ for some $x \in \mathcal{D}$. Then, by Theorem \ref{thm:c2-1-1}, there exists $r_0 > 0$ such that $A = u(x) = \mathcal{M}_{x,u}(r)$ for $0 < r \leq r_0$. But since $u$ is continuous and satisfies $u \leq A$, we must have $u \equiv A$ on $\Sigma_r(x) = \{y \in E_n; |y - x| = r\}$. That is, $u(t) = A$ for all $t = x + h$, with $|h| < r_0$. This shows that the set $\{t \in E_n; u(t) = A\}$ is open. On the other hand, the continuity of $u$ implies that this set is also closed (relative to $\mathcal{D}$). Since $\mathcal{D}$ is connected it follows that this set and $\mathcal{D}$ coincide. This shows that $u$ is the constant function $u(x) \equiv A$.

By applying this result to $-u$ we obtain the minimum principle for harmonic functions: Under the hypotheses of Corollary \ref{cor:c2-1-3}, if $B = \inf_{x\in\mathcal{D}} u(x) > -\infty$ then $u(x) > B$ for all $x\in\mathcal{D}$, provided $u$ is not a constant function.

An equivalent form of the maximum (minimum) principle is the following:

\begin{customcor}{1.3$'$}
\label{cor:c2-1-3p}
If u is continuous on the closure $\overline{\mathcal{D}}$ of the bounded domain $\mathcal{D}$ and is harmonic in $\mathcal{D}$ then the maximum (minimum) of u is attained only on the boundary $\partial\mathcal{D} = \overline{\mathcal{D}} - \mathcal{D}$ of $\mathcal{D}$, provided u is not a constant function.
\end{customcor}

Applying this to the difference $u_{1}-u_{2}$ we obtain:

\begin{corollary}
\label{cor:c2-1-4}
If $u_1$ and $u_2$ are continuous on the closure $\bar{\mathcal{D}}$ of the bounded domain $\mathcal{D}$, are harmonic in $\mathcal{D}$ and agree on the boundary $\partial\mathcal{D} = \bar{\mathcal{D}} - \mathcal{D}$. then $u_1(x) = u_2(x)$ for all $x \in \bar{\mathcal{D}}$.
\end{corollary}

Another consequence of the mean value property is the following result, known as Liouville's\index{Liouville theorem@Liouville theorem} theorem:

\begin{theorem}
\label{thm:c2-1-5}
If v is harmonic on $E_{n}$ and bounded then v is a constant function.
\end{theorem}

\textsc{Proof.}\ Let $\Omega_n = \omega_{n-1}/n$ be the volume of the solid unit sphere of $E_n$. From the mean value property Theorem \ref{thm:c2-1-1} we obtain, for all $x \in E_n$ and $t > 0$:

\[
\begin{aligned}v(x)&=v(x)\frac{n}{t^{n}}\int\limits_{0}^{t}\rho^{n-1}d\rho=\frac{n}{t^{n}}\int\limits_{0}^{t}\mathcal{M}_{x,v}(\rho)\rho^{n-1}d\rho\\&=\frac{n}{\omega_{n-1}}\frac{1}{t^{n}}\int\limits_{0}^{t}\left\{\int\limits_{\Sigma}v(x+\rho y^{\prime})dy^{\prime}\right\}\rho^{n-1}d\rho\\&=\frac{1}{t^{n}\Omega_{n}}\int\limits_{|y|\leq t}v(x+y)dy.\end{aligned}
\]

Thus, if $S_{t}(x)$ denotes the solid sphere about x of radius t,

\[
\begin{aligned}\left|v(x_{1})-v(x_{2})\right|&=\frac{1}{t^{n}\Omega_{n}}\Bigg|\int_{S_{t}(x_{1})}v(x)dx-\int_{S_{t}(x_{2})}v(x)dx\Bigg|\\&\leq\frac{1}{t^{n}\Omega_{n}}\Bigg|\Bigg(\int_{S_{t}(x_{1})-S_{t}(x_{2})}dx+\int_{S_{t}(x_{2})-S_{t}(x_{1})}dx\Bigg)\Bigg|\|v\|_\infty.\end{aligned}
\]

But the measure of the symmetric difference $(S_{t}(x_{1}) - S_{t}(x_{2})) \cup$

$(S_t(x_2) - S_t(x_1))$ divided by $t^n$ clearly tends to 0 as $t \to \infty$, and it follows that $v(x_1) - v(x_2) = 0$. Thus, v must be a constant function.

The mean value property completely characterizes harmonic functions. To see this we first establish the following lemma:

\begin{lemma}
\label{lem:c2-1-6}
Suppose $u$ is twice differentiable in the domain $\mathcal{D}$ and $\mathcal{M}_{x,u}(r) = u(x)$ whenever the sphere $\{t \in E_n; |t - x| \leq r\}$ belongs to $\mathcal{D}$, then $u$ is harmonic in $\mathcal{D}$.
\end{lemma}

\textsc{Proof.}\ Fix $x \in \mathcal{D}$ and let $\mathcal{M}''(0) (= \mathcal{M}_{x,u}''(0))$ be the limit, as $r \to 0$, of the second derivative $(d^2/dr^2)\mathcal{M}(r)$ (the continuity of the second derivatives of $u$ implies the existence of this limit). The lemma will then be an immediate consequence of the identity.

\[
\mathcal{M}(0)=\frac{1}{n}\left(\Delta u\right)(x).
\]

Letting $t' = (t_1', t_2', \ldots, t_n') \in \Sigma$ we have

\[
\frac{d^{2}}{d r^{2}}\mathcal{M}(r)=\frac{1}{\omega_{n-1}}\intop_{\Sigma}\left(\sum_{j,k=1}^{n}u_{j k}(x+r t^{\prime})t_{j}^{\prime}t_{k}^{\prime}\right)d t^{\prime}.
\]

Hence,

\[
\begin{aligned}{\mathcal{M}^{\prime\prime}(0)}&{{}=\frac{1}{\omega_{n-1}}\sum_{j,k=1}^{n}\int_{\Sigma}u_{j k}(x)t_{j}^{\prime}t_{k}^{\prime}\;{d}t^{\prime}}\\ {}&{{}=\frac{1}{\omega_{n-1}}\sum_{j,k=1}^{n}\left(\int_{\Sigma}t_{j}^{\prime}t_{k}^{\prime}\;{d}t^{\prime}\right)u_{j k}(x).}\\ \end{aligned}
\]

But, by an easy symmetry argument, $\int_{\Sigma} t_{j}^{\prime}t_{k}^{\prime}dt^{\prime}=0$ if $j\neq k$ while $\int_{\Sigma}(t_{j}^{\prime})^{2}dt^{\prime}=\omega_{n-1}/n$ is independent of $j$, $1\leq j\leq n$. Thus, we obtain $\mathcal{M}^{\prime\prime}(0)=(1/n)\sum_{j=1}^{n}u_{jj}(x)$ and the lemma is proved.

The following result shows, in particular, that Lemma \ref{lem:c2-1-6}. is valid even if we weaken the hypothesis of second differentiability by assuming only local integrability of $u$. Moreover, together with Theorem \ref{thm:c2-1-1}, the following result implies that a harmonic function on a domain $\mathcal{D}$ belongs to the class $C^{\infty}(\mathcal{D})$. We consider first the case when $u$ is continuous.

\begin{theorem}
\label{thm:c2-1-7}
Suppose a continuous function, $u$, on a domain $\mathcal{D} \subset E_n$ satisfies the mean value property (as in Lemma \ref{lem:c2-1-6}); then $u$ is harmonic and has continuous partial derivatives of all orders in $\mathcal{D}$.
\end{theorem}

\textsc{Proof.}\ Since the problem is local we can assume, by restricting u to a sphere whose closure belongs to $\mathcal{D}$, that $\mathcal{D}$ is a sphere and that $u$ is integrable on $\bar{\mathcal{D}}$. We can extend u to all of $E_n$ by defining it to be 0 outside

$\bar{\mathcal{D}}$. Choose a radial $\varphi\in C^\infty(E_n)$ such that $\int_{E_n}\varphi(x)\,dx=1$ and having support contained in the solid sphere of radius $1$. We then let $u_\varepsilon=u * \varphi_\varepsilon$, where $\varphi_\varepsilon(x)=\varepsilon^{-n}\varphi(x/\varepsilon)$ for $\varepsilon>0$. It is clear that $u_\varepsilon\in C^\infty(E_n)$. We have, using polar coordinates,

\[
\begin{aligned}{u_{\varepsilon}(x)=\int_{E_{n}}\varphi_{\varepsilon}(t)u(x-t)\;{d}t}&{{}=\int_{0}^{\varepsilon}\varphi_{\varepsilon}(r)\Biggl\{\int_{\Sigma}u(x-r t^{\prime})\;{d}t^{\prime}\Biggr\}r^{n-1}\;{d}r}\\ {}&{{}=\int_{0}^{\varepsilon}\varphi_{\varepsilon}(r)\omega_{n-1}\mathcal{M}_{x,u}(r)r^{n-1}\;{d}r.}\\ \end{aligned}
\]

If $\varepsilon$ is less than the distance from $x$ to the boundary of $\mathcal{D}$ then $\mathcal{M}_{x,i}(r) = u(x)$ for $0 \leq r \leq \varepsilon$. Thus, $u_{\varepsilon}(x) = u(x)\omega_{n-1}\int_{0}^{\varepsilon}\varphi_{\varepsilon}(r)r^{n-1}dr = u(x).\footnote{We already have seen in the first chapter (see Theorems \ref{thm:c1-1-18} and \ref{thm:c1-1-25}) how such convolutions $u_\varepsilon = u * \varphi_\varepsilon$ approximate $u$. With the added hypothesis that $u$ satisfies the mean value property we obtain the fact that $u_\varepsilon(x) = u(x)$ if $\varepsilon$ is small enough. In general, when $\varphi$ is smooth it follows that so is $u_\varepsilon = u * \varphi_\varepsilon$; in this case, such a family $\{u_\varepsilon\}$ is often called a regularization\index{regularization@regularization} of $u$.}$ In particular, if $2d$ is the radius and $x_0$ the center of $\mathcal{D}$, then $u(x) = u_d(x)$ for all $x$ whose distance from $x_0$ is less than $d$. This shows that in a neighborhood of each $x_0 \in \mathcal{D}$, $u$ coincides with a $C^\infty$ function, and the theorem is proved since the harmonicity of $u$ follows from Lemma \ref{lem:c2-1-6}.

If $u$ were merely locally integrable we would need to assume the mean-value property in the form

\[
u(x)=\frac{1}{t^{n}\Omega_{n}}\int\limits_{|y|<t}u(x+y)d y,
\]

for $t$ sufficiently small so that the sphere of center $x$ and radius $t$ lies in $\mathcal{D}$. By an argument analogous to the one given above this would show that $u(x) = u_{d}(x)$, and therefore $u$ would still satisfy the conclusions of Theorem \ref{thm:c2-1-7}.

\begin{corollary}
\label{cor:c2-1-8}
Suppose $\{u_n\}$ is a sequence of harmonic functions defined in the region $\mathcal{D}$. If this sequence converges uniformly to a function $u$ in each compact subset of $\mathcal{D}$ then $u$ is also harmonic.
\end{corollary}

\textsc{Proof.}\ Since $u$ is, then, continuous it follows from Theorem \ref{thm:c2-1-7} that it suffices to show that $u$ satisfies the mean value property. By Theorem \ref{thm:c2-1-1} $\mathcal{M}_{x,u_n}(r) = u_n(x)$ for all $x \in \mathcal{D}$ and $r$ such that $\{t; |x - t| \leq r\} \subset \mathcal{D}$. Since $u_n$ converges uniformly to $u$ on the sphere $\{t; |x - t| = r\}$, the

integrals defining the means $\mathcal{M}_{x,u_n}(r)$ must converge to the integral defining $\mathcal{M}_{x,u}(r)$. It follows that $\mathcal{M}_{x,u}(r) = \lim_{n \to \infty} \mathcal{M}_{x,u_n}(r) = \lim_{n \to \infty} u_n(x) = u(x)$, which is the desired result.

A classical problem involving harmonic functions is the well-known Dirichlet problem\index{Dirichlet problem@Dirichlet problem}. Its most common formulation is the following one:

Suppose $\mathcal{D}$ is a region with compact closure $\bar{\mathcal{D}}$ and $f$ a continuous function defined on the boundary $\partial\mathcal{D} = \bar{\mathcal{D}} - \mathcal{D}$. Does there exist a continuous function $u$ on $\bar{\mathcal{D}}$ such that

(i) u is harmonic when restricted to D.

(ii) $u(x) = f(x)$ when $x \in \partial \mathcal{D}$?

Corollary \ref{cor:c2-1-4} tells us that if such a function u exists then u is unique.

Our main interest lies in a variant of the Dirichlet problem, where $\mathcal{D}$ is the upper half-space $E_{n+1}^{+}$ (see Section 2). Nevertheless, the solution of the problem as stated here, when $\mathcal{D}$ is the interior of a sphere\index{Poisson kernel@Poisson kernel!for sphere@for sphere}, is of basic importance and has many applications. In this special case, we obtain the solution by making use of the Poisson kernel for the unit sphere:

\[
p(s,x)=\frac{1}{\omega_{n-1}}\frac{1-|x|^{2}}{|x-s|^{n}}=\frac{1}{\omega_{n-1}}\frac{1-r^{2}}{(1-2r\cos\gamma+r^{2})^{n/2}},
\]

where $r = |x| < 1 = |s|$ and $\gamma$ is the angle between the vectors x and s; that is, $r \cos \gamma = x \cdot s$.

Three basic properties of this kernel are the following:

\begin{theorem}
\label{thm:c2-1-9}
If $\Sigma = \Sigma_{n-1} = \{s \in E_n; |s| = 1\}$ denotes the surface of the unit sphere of $E_n$ and ds the element of surface area, then
(a) $p(s, x) \geq 0$ for all $s \in \Sigma$ and $|x| < 1$;
(b) $\int_{\Sigma} p(s, x) \, ds = 1$ whenever $|x| < 1$;
(c) for $x' \in \Sigma$ let $x = rx' (0 \leq r < 1)$ then, for each $\delta > 0$,
\[
\int\limits_{s\in\Sigma,|s-x^{\prime}|>\delta}p(s,r x^{\prime}) d s\to0
\]

as $r \to 1$, uniformly in $x'$
\end{theorem}

\textsc{Proof.}\ Property (a) is obvious and so is (c) once we observe that $|s - x| = |s - rx'|$ is bounded away from 0 (uniformly in $x'$) when $s \in \Sigma$ satisfies $|s - x'| > \delta$.

In order to show (b) we first note that $p(s, x)$ is a harmonic function of $x, |x| < 1$, for each $s \in \Sigma$. Thus, by the mean value property (Theorem \ref{thm:c2-1-1}),

\[
1\;=\;\omega_{n-1}p(s,0)=\frac{1}{\omega_{n-1}}\intop_{\Sigma}\omega_{n-1}p(s,r x^{\prime})\;d x^{\prime}\;=\;\intop_{\Sigma}p(s,r x^{\prime})\;d x^{\prime}
\]

for $0 \leq r < 1$ and $s$, $x' \in \Sigma$. But $|rx' - s| = |rs - x'|$ and, consequently, $p(s, rx') = p(x', rs)$. We obtain (b), therefore, by interchanging the roles of $s$ and $x'$.

Using these properties we can now solve the Dirichlet problem for the interior of the unit sphere:

\begin{theorem}
\label{thm:c2-1-10}
Suppose f is a continuous function on $\Sigma$, then the function $u$ defined by
\[
u(x)=\int\limits_{\Sigma}f(s)p(s,x)\,ds
\]

when $|x| < 1$ and $u(x) = f(x)$ when $|x| = 1$ is harmonic for $|x| < 1$ and continuous for $|x| \leq 1$.
\end{theorem}

\textsc{Proof.}\ The fact that $u$ is harmonic for $|x|<1$ follows from the fact that the Laplacian of $u$ is equal to the integral of the product of $f$ with the Laplacian (with respect to $x$) of $p(s,x)$. This requires justifying the change in order of differentiation and integration. A simpler proof is obtained by observing that it follows immediately from Fubini's theorem that the mean value of $u$ taken over a sphere of radius $r<1-|x|$ about the point $x$ equals the integral of $f(s)$ times the mean value of $p(s,\cdot)$ over the same sphere. Since $p(s,\cdot)$ is harmonic the latter equals $p(s,x)$. Thus, $u$ satisfies the mean value property and, by Theorem \ref{thm:c2-1-7}, must, therefore, be harmonic in the interior of the unit sphere.

To show that $u$ is continuous on the closed unit sphere it suffices to show that this is true on the boundary $\Sigma$. Let $x' \in \Sigma$, $0 \leq r < 1$, $x = rx'$ and $A_\delta = \{s \in \Sigma; |s - x'| > \delta\}$. Then, by (b) and (a) of Theorem \ref{thm:c2-1-9}:

\[
\begin{aligned}|u(x)-u(x^{\prime})|&=\left|\int\limits_{\Sigma}[f(s)-f(x^{\prime})]p(s,rx^{\prime})ds\right|\\&\leq\int\limits_{A_\delta}|f(s)-f(x^{\prime})|p(s,rx^{\prime})ds\\&\quad+\int\limits_{\Sigma-A_\delta}|f(s)-f(x^{\prime})|p(s,rx^{\prime})ds.\end{aligned}
\]

Again making use of (b), Theorem \ref{thm:c2-1-9}, we see that the second summand is majorized by

\[
\left\{\sup_{|s-x^{\prime}|\leq\delta}|f(s)-f(x^{\prime})|\right\}\intop_{\Sigma-A_{\delta}}p(s,x)\mathop{}{d}s\leq\left\{\sup_{|s-x^{\prime}|\leq\delta}|f(s)-f(x^{\prime})|\right\}\cdot1,
\]

which is small, say less than $\varepsilon > 0$, if $\delta$ is chosen close to 0. On the other hand, for $\delta$ fixed, the first summand is majorized by

\[
2\biggl\{\sup_{t\in\Sigma}|f(t)|\biggr\}\intop_{A_{\delta}}p(s,x)\mathop{}{d}s;
\]

but, by Theorem \ref{thm:c2-1-9} part (c), the latter tends to 0 as $r \to 1$, uniformly in $x'$. Thus, $|u(x) - u(x')| < \varepsilon$, uniformly in $x'$, if r is close enough to 1.

Consequently, if $x_0 \in \Sigma$ and $x = rx'$ we then have $|u(x_0) - u(x)| \leq |u(x_0) - u(x')| + |u(x') - u(x)| = |f(x_0) - f(x')| + |u(x') - u(x)| < \varepsilon + \varepsilon = 2\varepsilon$ if $x$ is sufficiently close to $x_0$ (since this implies that both 1 - $r$ and $|x' - x|$ are small). This proves the theorem.

By applying an appropriate translation and dilation, this result gives us the solution of the Dirichlet problem for the general sphere in $E_{n}$:

\begin{corollary}
\label{cor:c2-1-11}
Let $\mathcal{D}$ be the interior of the sphere of radius $a$ about $x_0$ and $f$ a continuous function on the boundary, $\partial\mathcal{D}$, of $\mathcal{D}$. Then the function $u$ defined by
\[
u(x)=\frac{1}{\omega_{n-1}a^{2-n}}\;\int\limits_{\Sigma}f(x_{0}+a s)\;\frac{a^{2}-|x-x_{0}|^{2}}{|(x-x_{0})-a s|^{n}}d s
\]

if $|x - x_0| < a$ and $u(x) = f(x)$ if $|x - x_0| = a$, is harmonic in $\mathcal{D}$ and continuous on $\mathcal{D}$.
\end{corollary}

We shall give two applications of this result. The first is the following basic theorem concerning uniformly bounded sequences of harmonic functions:

\begin{theorem}
\label{thm:c2-1-12}
Suppose $\{u_m\}$ is a sequence of harmonic functions defined in a domain $\mathcal{D} \subset E_n$ such that the functions $u_m$ are uniformly bounded in the closure of a bounded subdomain $\mathcal{S}$, such that $\bar{\mathcal{S}} \subset \mathcal{D}$. Then there exists a subsequence $\{u_{m_k}\}$ that converges uniformly to a harmonic function defined on $\mathcal{S}$.
\end{theorem}

\textsc{Proof.}\ If we can show that the functions $u_m$ are equicontinuous on $\mathcal{F}$ then this theorem is a consequence of Ascoli's theorem and Corollary \ref{cor:c2-1-8}. In order to show this equicontinuity it suffices to show that the first derivatives of the $u_m$ are uniformly bounded in each closed subset of $\mathcal{S}$. But, if $x_0 \in \mathcal{S}$ then $u_m(x)$ has the integral expression of Corollary \ref{cor:c2-1-11}, for $|x - x_0| < a$, as long as $a$ is small enough so that the sphere of radius $a$ about $x_0$ lies in $\mathcal{S}$. Differentiating with respect to $x_i$ ($i = 1, 2, \ldots, n$) and evaluating this partial derivative at $x_0$, we obtain

\[
\frac{\partial u_{m}}{\partial x_{i}}\left(x_{0}\right)=\frac{n}{\omega_{n-1}a}\int\limits_{\Sigma}u_{m}(x_{0}+a s)s_{i}d s.
\]

Thus, $|(\partial u_m/\partial x_i)(x_0)| \leq (n/a) \sup_{y \in \mathcal{S}} |u_m(y)|$ from which we obtain the desired uniform boundedness of the first partial derivatives of the $u_m$'s in each closed subset of $\mathcal{S}$.

The second application of Corollary \ref{cor:c2-1-11} is the following reflection principle\index{reflection principle@reflection principle} for harmonic functions:

\begin{theorem}
\label{thm:c2-1-13}
Suppose $\mathcal{D} \subset E_{n+1}$ is a domain that is symmetric with respect to $E_n$ (that is, if $(x, y) = (x_1, x_2, \ldots, x_n, y) \in \mathcal{D}$ then $(x, -y) \in \mathcal{D}$). If a continuous function, $u$, defined on $\mathcal{D}$ satisfies $u(x, y) = -u(x, -y)$ and is harmonic in the ``upper half,'' $\mathcal{D}^+ = \{(x, y) \in \mathcal{D}; y > 0\}$ of $\mathcal{D}$ then it is harmonic in all of $\mathcal{D}$.
\end{theorem}

\textsc{Proof.}\ The function $u$ is harmonic in $\mathcal{D}^+$ by assumption. Since $-(\Delta u)(x, -y) = (\Delta u)(x, y)$, it is harmonic in $\mathcal{D}^-$ = $\{(x, y) \in \mathcal{D}; y < 0\}$ as well. Consequently, it suffices to show that $u$ is harmonic in a neighborhood of each point of $\mathcal{D} \cap E_n = \{(x, y) \in \mathcal{D}; y = 0\}$. If $x_0 = (x_0, 0)$ is such a point then we can find a closed (solid) sphere, $S_a(x_0)$, of radius $a > 0$ centered at this point that is contained in $\mathcal{D}$. For $(x, y)$ in this sphere let

\begin{equation}\label{eq:c2-1-14}
w(x,y)=\frac{a^{n-1}}{\omega_{n}}\intop_{\Sigma_{n}}u(x_{0}+a s,a t)\frac{a^{2}-|(x-x_{0},y)|^{2}}{|(x-x_{0}-a s,y}-a t)|^{n+1}d\sigma,
\end{equation}

where $\sigma = (s, t) = (s_1, s_2, \ldots, s_n, t)$ is a point on the surface, $\Sigma_n$, of the unit sphere of $E_{n+1}$. By Corollary \ref{cor:c2-1-11}, $w$ is the harmonic function agreeing with $u$ on the boundary of $S_a(x_0)$. On the other hand, the integral \eqref{eq:c2-1-14} vanishes when $y = 0$, since $u(x_0 + as, at) = -u(x_0 + as, -at)$,

\[
\frac{a^{2}-|(x-x_{0},0)|^{2}}{|(x-x_{0}-as,-at)|^{n+1}}=\frac{a^{2}-|(x-x_{0},0)|^{2}}{|(x-x_{0}-as,at)|^{n+1}},
\]

while by continuity, $u(x_0,0)=0$ whenever $(x_0,0)\in\mathcal{D}$. Thus, $w$ agrees with the continuous function $u$ on the upper hemisphere $\{(x,y)\in\bar{E}_{n+1}^+ ; |x-x_0|^2+y^2=a^2\}$, consisting of the upper half of the boundary of $S_a(x_0)$, and on $S_a(x_0)\cap E_n$. Since both $w$ and $u$ are harmonic in the interior of this hemisphere, by Corollary \ref{cor:c2-1-4} they must agree there. Applying the same argument to the lower hemisphere we obtain $u\equiv w$ on $S_a(x_0)$. Since $w$ is harmonic in the interior of $S_a(x_0)$, so is $u$ and the theorem is proved.

\begin{corollary}
\label{cor:c2-1-15}
Suppose u is a function that is continuous on $\overline{E}_{n+1}^{+} = E_{n+1}^{+} \cup E_n = \{(x, y) \in E_{n+1}; y \geq 0\}$ that is harmonic in $E_{n+1}^{+}$ and vanishes on $E_n$. Then, if u is bounded in $E_{n+1}^{+}$ it follows that u is identically 0.
\end{corollary}

Let us note that unless a restriction, such as boundedness, is imposed on

u this result does not hold. For example, the function $u(x, y) = y$ satisfies all the conditions of the corollary except that it is not bounded. This example also shows that the Dirichlet problem, posed in the case when $\mathcal{D}$ is the unbounded domain $E_{n+1}^+$, does not have a unique solution (for $u(x, y) = y$ and $v(x, y) = 0$ are harmonic in $\mathcal{D}$ and both have boundary value $f \equiv 0$ on $E_n$). Corollary \ref{cor:c2-1-15}, on the other hand, tells us that the solution is unique if we impose on it the added condition that it be bounded. In the next section we make a detailed study of some variants of the Dirichlet problem associated with the region $E_{n+1}^+$.

Now that we have disposed of the required preliminaries we come to that part of the theory of harmonic functions connected with Fourier analysis on $E_{n}$.

\section{The Characterization of Poisson Integrals}

We pose the following problem: Given a function $f \in L^p(E_n)$, $1 \leq p \leq \infty$, does there exist a harmonic function, $u$, defined on the upper half space $E_{n+1}^+$ such that

\[
\left\|u(\cdot,y)-f\right\|_{p}=\left(\int_{E_{n}}\left|u(x,y)-f(x)\right|^{p}d x\right)^{1/p}\to0
\]

as $y \to \mathbb{C}$? If $f$ were continuous and $p = \infty$ then this would be a direct generalization of the Dirichlet problem posed for the unbounded domain $\mathcal{D} = E_{n^+_1}$. We have just seen that the solution need not be unique in this case. It is also clear that if $f$ is not continuous then there could not exist a harmonic (and, in particular, continuous) function $u$ such that $\|u(\cdot, y) - f\|_\infty \to 0$ as $y \to 0$. Nevertheless, a solution of this problem does exist when $p < \infty$; if the problem is properly posed, a solution also exists if $p = \infty$.

The following result, in large part a collection of results that have been already established, tells us that the Poisson integral of f is the harmonic function we are looking for.

\begin{theorem}
\label{thm:c2-2-1}
(a) If $f \in L^p(E_n)$, $1 \leq p \leq \infty$, and
\[
u(x,y)=\int\limits_{E_{n}}f(t)P(x-t,y)dt
\]

is the Poisson integral of f then u is harmonic in $E_{n+1}^{+}$, $\lim_{y\to0}u(x,y)=f(x)$ for almost every $x\in E_{n}$ and
\begin{equation}\label{eq:c2-2-2}
\left\|u(\cdot,y)\right\|_{p}=\left(\int_{E_{n}}\left|u(x,y)\right|^{p} dx\right)^{1/p}\leq\left\|f\right\|_{p}
\end{equation}

for all $y > 0$. If $1 \leq p < \infty$ then $u(\cdot, y)$ converges to $f$ in the $L^p$ norm as
$y \to 0$; that is,
\[
\|u(\cdot,y)-f\|_{p}=\left(\int_{E_{n}}|u(x,y)-f(x)|^{p} dx\right)^{1/p}\to0
\]

as $y \rightarrow 0$.
(b) If $f \in C_0 \subset L^\infty(E_n)$ then the Poisson integral of $f$, $u(x, y)$, converges to $f$ uniformly: $\|u(\cdot, y) - f\|_\infty = \sup_{x \in E_n} |u(x, y) - f(x)| \to 0$ as $y \to 0$. If we only assume $f$ to be continuous and bounded the convergence is uniform on compact subsets of $E_n$. In either case, we can extend $u(x, y)$ to $\overline{E_{n+1}^+} = E_{n+1}^+ \cup E_n$ by letting $u(x, 0) = f(x)$ and we obtain a continuous function.
\end{theorem}

\textsc{Proof.}\ The almost everywhere convergence and the convergence in the norm in part (a) are Theorems \ref{thm:c1-1-25} and 1.18 in Chapter I. Inequality \eqref{eq:c2-2-2} is an immediate consequence of Theorem \ref{thm:c1-1-3} in that chapter, with $g(x) = P(x, y)$ and $h(x) = u(x, y)$ (recall that Lemma \ref{lem:c1-1-17} asserts that $g = P(\cdot, y) \in L^1(E_n)$ and $\|g\|_1 = 1$). The fact that u is harmonic can be shown by using the harmonicity of the Poisson kernel and an argument completely analogous to that used in the beginning of the proof of Theorem \ref{thm:c2-1-10}.

If $f$ is continuous and bounded $u(x, y)$ converges to $f(x)$, as $y \to 0$, for all $x \in E_n$, since condition \eqref{eq:c1-1-24} of Chapter I obviously holds everywhere (see Theorem \ref{thm:c1-1-25} in that chapter). The uniform convergence on a compact subset, $F$, of $E_n$ needs only the following simple argument: given $\varepsilon > 0$ there exists $\delta > 0$ such that $|f(x - t) - f(x)| < \varepsilon$ whenever $x \in F$ and $|t| < \delta$. Thus,

\[
\begin{aligned}|u(x,y)-f(x)|&=\left|\int\limits_{E_{n}}f(x-t)P(t,y)dt-f(x)\right|\\&=\left|\int\limits_{E_{1}}f(x-t)P(t,y)dt-\int\limits_{E_{n}}f(x)P(t,y)dt\right|\\&\leq\int\limits_{|t|\leq\delta}|f(x-t)-f(x)|P(t,y)dt\\&\quad+\int\limits_{|t|>\delta}|f(x-t)-f(x)|P(t,y)dt\\&\leq\varepsilon\int\limits_{E_{n}}P(t,y)dt+2\|f\|_{\infty}\int\limits_{|t|>\delta}P(t,y)dt\\&=\varepsilon\cdot1+2\|f\|_{\infty}\int\limits_{|t|>\delta}P(t,y)dt.\end{aligned}
\]

But, for $\delta > 0$ fixed, $\int_{|t| > \delta} P(t, y) \, dt \leq c_n y \int_{|t| > \delta} |t|^{-(n+1)} \, dt$ tends to 0 as $y \to 0$. Consequently, $|u(x, y) - f(x)| \leq \varepsilon$, if $y$ is sufficiently small, for all $x \in F$. If $f \in C_0$ then it is uniformly continuous and, consequently, we can find $\delta > 0$ such that $|f(x-t) - f(x)| < \varepsilon$ for all $x \in E_n$ when $|t| < \delta$; we therefore obtain the uniform convergence of $u(x, y)$ to $f(x)$ as $y \to 0$. This proves the theorem.\footnote{The convergence $\|u(\cdot, y) - f\|_\infty \to 0$ when $f \in C_0$ is a special case of Theorem \ref{thm:c1-1-18} of Chapter I. The proof we have given of part (b) was based on the following three properties of the Poisson kernel: (i) $P(x, y) \geq 0$; (ii) $\int_{E_n} P(x, y) \, dx = 1$ for all $y > 0$; (iii) if $\delta > 0$ then $\int_{|x| > \delta} P(x, y) \, dx \to 0$ as $y \to 0$ (i) is obvious, (iii) was just proved and (ii) is part (b) of Lemma \ref{lem:c1-1-17} of Chapter I). A kernel satisfying these three properties is often called an approximation to the identity. It is clear that the Gauss--Weierstrass kernel also satisfies these properties. The reader will note that a slight modification of the last part of the proof of part (b) (replacing the $L^\infty$ norm by the $L^p$ norm, $1 \leq p < \infty$) gives us another proof of Theorem \ref{thm:c1-1-18} when $\varphi(x) = P(x, 1)$. An advantage of this treatment is that it does not make use of the operation of convolution (see, for example, Theorem \ref{thm:c2-1-9}, where analogous properties (a), (b), and (c) enable us to obtain the solution of the Dirichlet problem for the sphere: Theorem \ref{thm:c2-1-10}). Later, in Theorem \ref{thm:c3-5-6} of Chapter III this type of $L^p$ convergence is established by these methods for another kind of Poisson integral.}

The space $L^{1}(E_{n})$ is naturally embedded in the space $M(E_{n})$ of finite Borel measures on $E_{n}$ (see the proof of Theorem \ref{thm:c2-3-19} in the last chapter). Many results and concepts associated with the former space can be extended to the\index{approximation to the identity@approximation to the identity} latter. In the first chapter for example, we have done so with the Fourier transform, the operation of convolution and Theorem \ref{cor:c2-1-3}. Similarly, we can do this for Theorem \ref{thm:c2-2-1}:

\begin{theorem}
\label{thm:c2-2-3}
If $\mu \in M(E_n)$ and $u(x, y) = \int_{E_n} P(x - t, y) \, d\mu(t)$ is the Poisson--Stieltjes integral of $\mu$ then $u$ is harmonic in $E_{n+1}^{+}$ and
\begin{equation}\label{eq:c2-2-4}
\|u(\cdot,y)\|_{1}=\int\limits_{E_{n}}|u(x,y)|d x\leq\|\mu\|,
\end{equation}

where $\|\mu\|$ is the total variation, $|\mu|(E_{n})$, of $\mu$. Furthermore,
\[
\lim_{y\to0}\int\limits_{E_{n}}u(x,y)\varphi(x)dx=\int\limits_{E_{n}}\varphi(x)d\mu(x)
\]

for all $\varphi \in C_0$; that is, $u(\cdot, y)$ converges to $\mu$ in the weak$\cdot$ topology.
\end{theorem}

\textsc{Proof.}\ The fact that $u$ is harmonic can be shown by the same argument that was used in the proof of Theorem \ref{thm:c2-1-10} (and of Theorem \ref{thm:c2-2-1}). Inequality \eqref{eq:c2-2-4} is an immediate consequence of the extension to Borel measures of Theorem \ref{thm:c1-1-3} (in Chapter I) that was described after the

statement of that theorem. The weak* convergence can be shown in the following way: If $\varphi \in C_0$ and $v(x, y)$ is its Poisson integral then

\[
\begin{aligned}\int_{E_{n}}u(x,y)\varphi(x)dx=&\int_{E_{n}}\left(\int_{E_{n}}P(x-t,y)d\mu(t)\right)\varphi(x)dx\\=&\int_{E_{n}}\left(\int_{E_{n}}P(x-t,y)\varphi(x)dx\right)d\mu(t)\\=&\int_{E_{n}}v(t,y)d\mu(t).\end{aligned}
\]

Consequently, since, by part (b) of Theorem \ref{thm:c2-2-1}, $\|v(\cdot, y) - \varphi\|_\infty \to 0$ as $y \to 0$,

\[
\begin{aligned}\left|\int_{E_{n}}u(x,y)\varphi(x)dx-\int_{E_{n}}\varphi(t)d\mu(t)\right|&=\left|\int_{E_{n}}\{v(t,y)-\varphi(t)\}d\mu(t)\right|\\&\leq\int_{E_{n}}|v(t,y)-\varphi(t)||d\mu|(t)\\&\leq\|v(\cdot,y)-\varphi\|_{\infty}|\mu|(E_{n})\to0\end{aligned}
\]

as $y \to 0$ and the theorem is established.

The principal result of this section is the proof of a converse of Theorem \ref{thm:c2-2-1}. Inequality \eqref{eq:c2-2-2} asserts that the $L^p$ norms $\|u(\cdot, y)\|_p$ are bounded for $y > 0$. We show that this is sufficient for $u$ to be a Poisson integral.

\begin{theorem}
\label{thm:c2-2-5}
If $u(x, y)$ is harmonic in $E_{n+1}^+$ and there exists a constant $c > 0$ and $a p, 1 \leq p \leq \infty$, such that
\[
\|u(\cdot,y)\|_{p}=\left(\int_{E_{n}}|u(x,y)|^{p}d x\right)^{1/p}\leq c<\infty
\]

for all y > 0 then
(a) when $1 < p \leq \infty$, $u(x, y)$ is the Poisson integral of a function $f$ in $L^p(E_n)$;
(b) if $p = 1$, $u(x, y)$ is the Poisson--Stieltjes integral of a finite Borel measure; if, in addition, $u(\cdot, y)$ is Cauchy in the $L^1$ norm as $y \to 0$ then $u(x, y)$ is the Poisson integral of a function $f$ in $L^1(E_n)$.
\end{theorem}

\textsc{Proof.}\ If $1 < p \leq \infty$ and $\|u(\cdot, y)\|_p \leq c < \infty$ for all $y > 0$ then there exists a sequence $\{y_k\}$, with $\lim_{k \to \infty} y_k = 0$, and a function $f$ in $L^p(E_n)$ such that $u(\cdot, y_k)$ converges weakly to $f$ as $k \to \infty$. That is, for each $g \in L^{p'}(E_n)$,

\[
1/p\;+\;1/p^{\prime}=1,
\]

\[
\lim_{k\to\infty}\int\limits_{E_{n}}u(x,y_{k})g(x)dx=\int\limits_{E_{n}}f(x)g(x)dx.
\]

If $p = 1$ there exists a finite Borel measure $\mu$ that is the weak* limit of a sequence $\{u(\cdot, y_k)\}$. That is, for each $g$ in $C_0$

\[
\lim_{k\to\infty}\int\limits_{E_{n}}u(x,y_{k})g(x)dx=\int\limits_{E_{n}}g(x)d\mu(x).\footnote{These results are special cases of the fact that the unit sphere of the dual of a Banach space is compact in the weak* topology. The last assertion, concerning p = 1, is often known as Helly's\index{Helly theorem@Helly theorem} theorem (see bibliographical notes to Chapter I).}
\]

Since $P(\cdot, y)$, for each $y > 0$, belongs to $L^{p'}(E_n)$ for $1 \leq p' \leq \infty$ and also belongs to $C_0$ we have, in particular,

\[
\lim_{k\to\infty}\int\limits_{E_{n}}P(x-t,y)u(t,y_{k})dt=\int\limits_{E_{n}}P(x-t,y)f(t)dt=v(x,y)
\]

when $1 < p \leq \infty$, and

\[
\lim_{k\to\infty}\int\limits_{E_{n}}P(x-t,y)u(t,y_{k})dt=\int\limits_{E_{n}}P(x-t,y)d\mu(t)=v(x,y)
\]

when p = 1.

If we show that, in either case, $v(x, y) = u(x, y)$ then we will have established (a), and the first part of (b). But this is an immediate consequence of the following two lemmas (since they imply that $u(x, y + y_k) = \int_{E_n} P(x - t, y) u(t, y_k) \, dt)$:

\begin{lemma}
\label{lem:c2-2-6}
If $u(x, y)$ satisfies the hypotheses of Theorem \ref{thm:c2-2-5} then there exists a constant $A = A_{n,p} > 0$ such that
\[
\left\|u(\cdot,y)\right\|_{\infty}=\sup_{x\in E_{n}}\left|u(x,y)\right|\leq A c y^{-n/p}.
\]

\end{lemma}

In particular, u is bounded in each proper sub-half-space

\[
E_{n+1,y_{0}}^{+}=\{(x,y)\in E_{n+1};y\geq y_{0}>0\}\subset E_{n+1}^{+}.
\]

\begin{lemma}
\label{lem:c2-2-7}
If $u(x, y)$ is harmonic in $E_{n+1}^{+}$ and bounded in each proper sub-half-space then
\[
u(x,y_{1}+y_{2})=\int\limits_{E_{n}}u(t,y_{1})P(x-t,y_{2})dt
\]

for all $y_{1}, y_{2} > 0$.
\end{lemma}

The first of these lemmas is an easy consequence of the mean value property for harmonic functions: Let $\Omega_{n+1}$ be the volume of the solid unit sphere of $E_{n+1}$, then $\Omega_{n+1} = \omega_n/(n+1)$ and

\[
\begin{aligned}u(x,y)&=u(x,y)\frac{n+1}{(y/2)^{n+1}}\int\limits_{0}^{y/2}r^{n} dr\\&=\left[(n+1)2^{n+1}/y^{n+1}\right]\int\limits_{0}^{y/2}\mathcal{M}_{(x,y),u}(r)r^{n} dr\\&=\left[(n+1)2^{n+1}/y^{n+1}\omega_{n}\right]\\&\quad\times\int\limits_{0}^{y/2}\left\{\int\limits_{\Sigma_{n}}\boldsymbol{u}(x_{1}+rt_{1}^{\prime},\ldots,x_{n}+rt_{n}^{\prime},y+rt_{n+1}^{\prime})dt^{\prime}\right\}r^{n} dr\\&=\left(2^{n+1}/(\Omega_{n+1})y^{-(n+1)}\int\limits_{[0]\leq y/2}u(x_{1}+t_{1},\ldots,x_{n}+t_{n},y+t_{n+1})dt.\right.\end{aligned}
\]

Consequently, letting $(\xi_1, \ldots, \xi_n, \eta) = (\xi, \eta)$ and $A = (2^{n+1}/\Omega_{n+1})^{1/p}$,

\[
\begin{aligned}|u(x,y)|&\leq(2^{n+1}/\Omega_{n+1})y^{-(n+1)}\int\limits_{[(x,y)-(\xi,\eta)]<y/2}|u(\xi,\eta)|d\xi d\eta\\&\leq(2^{n+1}/\Omega_{n+1})y^{-(n+1)}\int\limits_{[{(x,y)-(\xi,\eta)}<y/2}|u(\xi,\eta)|^{p} d\xi d\eta\Bigg)^{1/p}\\&\quad\times((\Omega_{n+1}y^{n+1}/2^{n+1})^{1-(1/p)}\\&\leq A y^{-(n+1)/p}\Bigg[\int\limits_{y/2}^{3y/2}\left\{\int\limits_{\xi\in E_{n}}|u(\xi,\eta)|^{p} d\xi\right\}d\eta\Bigg]^{1/p}\\&\leq A y^{-(n+1)/p}\Bigg[\int\limits_{y/2}^{3y/2}c^{p} d\eta\Bigg]^{1/p}=A c y^{-(n+1)/p}y^{1/p}=A c y^{-n/p}.\end{aligned}
\]

In order to prove \eqref{lem:c2-2-7} we fix $y_0 > 0$ and put $w(x, y) = u(x, y + y_0)$ for $y \geq 0$. Let $w_1(x, y) = \int_{E_n} u(t, y_0) P(x - t, y) \, dt$ for $(x, y) \in E_{n+1}^+$). The lemma will be established if we show that $w \equiv w_1$.

It follows from Lemma \ref{lem:c2-2-6} that the function $w$ is harmonic in $E_{n+1}^{+}$ and, on $\overline{E_{n+1}^{+}}$, it is continuous and bounded. From Theorem \ref{thm:c2-2-1} (see parts (a) and (b)) $w_{1}$ can be extended to $\overline{E_{n+1}^{+}}$ by letting $w_{1}(x,0)=u(x,y_{0})$ and we obtain a continuous and bounded function that is harmonic in $E_{n+1}^{+}$. Thus, the function $h=w_{1}-w$ is also continuous and bounded on $\overline{E_{n+1}^{+}}$ and harmonic in $E_{n+1}^{+}$. Since $h(x,0)=0$ for all $x\in E_{n}$ it follows from Corollary \ref{cor:c2-1-15} that $h$ is identically zero; thus, we obtain the desired result $w\equiv w_{1}$.

We still need to show the last part of Theorem \ref{thm:c2-2-5}. If $u(\cdot, y)$ is Cauchy in the $L^1$ norm it follows from the completeness of $L^1(E_n)$ that there exists $f \in L^1(E_n)$ such that $\|u(\cdot, y) - f\|_1 \to 0$ as $y \to 0$. Therefore, $\int_{E_n} u(x, y) g(x) \, dx \to \int_{E_n} f(x) g(x) \, dx$ whenever $g \in L^\infty(E_n)$. Now, the same argument that was used to establish part (a) of Theorem \ref{thm:c2-2-5} applies and we can conclude that $u$ is the Poisson integral of the $L^1$ function $f$.

\section{The Hardy-Littlewood Maximal Function and Nontangential Convergence of Harmonic Functions}\index{Hardy-Littlewood maximal function@Hardy-Littlewood maximal function!mf@$m_f$}

Up to this point we have considered only one type of convergence at the boundary for harmonic functions defined in $E_{n+1}^{+}$. We have shown, for example, that when u is a Poisson integral then, for almost every point $x_0 \in E_n$, there exists a boundary value

\begin{equation}\label{eq:c2-3-1}
\lim\limits_{(x,y)\to(x_{0},0)}u(x,y)
\end{equation}

as the point $(x, y)$ approaches $(x_0, 0)$ ``downward'' along the line $x = x_0$ (see Theorem \ref{thm:c2-2-1}). An important more general result holds, however; the limits \eqref{eq:c2-3-1} exist almost everywhere even when $(x, y)$ is allowed to approach $(x_0, 0)$ along more general paths. These are the nontangential paths. Before making this notion more precise we shall study an important operator, the Hardy--Littlewood maximal operator, whose basic properties can be used to show the existence of these more general limits.

This operator assigns to each function $f \in L^p(E_n)$, $1 \leq p \leq \infty$, its Hardy-Littlewood maximal function $m_f$ defined at $x \in E_n$ by

\[
\begin{aligned}m_{j}(x)&=\sup_{\tau>0}\frac{1}{\Omega_{n} r^{n}}\int\limits_{|t|\leq r}\left|f(x-t)\right|dt\\&=\sup_{S_{x}}\frac{1}{\left|S_{x}\right|}\int\limits_{S_{x}}\left|f(t)\right|dt,\end{aligned}
\]

where $\Omega_n$ is the ``volume'' (Lebesgue measure) of the solid unit sphere $\{t \in E_n: |t| \leq 1\}$; the last supremum is taken over all spheres $S_x$ having center $x$ and positive radius, and $|S_x|$ is the Lebesgue measure of $S_x$ (thus, if the radius of $S_x$ is $r$ we have $|S_x| = \Omega_n r^n$). We sometimes refer to $m_f$ as the spherical maximal function in order to distinguish it from the cubic maximal function $\bar{m}_f$ defined by

\[
\overline{{m}}_{f}(x)=\sup_{Q_{x}}\frac{1}{\left|Q_{x}\right|}\intop_{Q_{x}}\left|f(t)\right|dt,
\]

where the supremum is taken over all nondegenerate cubes $Q_x$ centered at $x$ having sides parallel to the axes and $|Q_x|$ denotes the Lebesgue measure

of $Q_x$. It is clear that, if $f$ is essentially bounded, then $m_f(x)$ and $\overline{m}_f(x)$ do not exceed $\|f\|_\infty$. We shall show that $m_f(x)$ and $\overline{m}_f(x)$ are finite almost everywhere whenever $f\in L^p(E_n)$, $1\leq p\leq\infty$.

Given a sphere $S_x$ of radius $r > 0$ and center $x$, let $q_x$ be the inscribed cube (it has, therefore, diagonals of length $2r$) and $Q_x$ the circumscribed cube (it has, therefore, sides of length $2r$), both cubes having sides parallel to the axes. Then, clearly, there exist constants $a_n$ and $A_n$, depending only on the dimension, such that

\[
|Q_{x}|\leq A_{n}|S_{x}|\quad\text{and}\quad|S_{x}|\leq a_{n}|q_{x}|.
\]

Thus,

\[
\frac{1}{\left|S_{x}\right|}\int\limits_{S_{x}}\left|f(t)\right|dt\leq\frac{A_{n}}{\left|Q_{x}\right|}\int\limits_{Q_{x}}\left|f(t)\right|dt.
\]

But this implies

\begin{equation}\label{eq:c2-3-2}
m_{f}(x)\leq A_{n}\overline{{m}}_{f}(x).
\end{equation}

Similarly, we have

\begin{equation}
\tag{3.2$'$}\label{eq:c2-3-2p}
\overline{m}_{f}(x)\leq a_{n}m_{f}(x).
\end{equation}

We shall be more interested in the spherical maximal function $m_f$. In some cases, however, it is easier to establish needed estimates for $\overline{m}_f$; inequalities \eqref{eq:c2-3-2} and \eqref{eq:c2-3-2p} can then be used to obtain corresponding results for $m_f$. The following lemma will enable us to show that $\overline{m}_f(x) < \infty$ almost everywhere (and, hence, $m_f(x) < \infty$ a.e.).

\begin{lemma}
\label{lem:c2-3-3}
Let $R_1 \subset R_2 \subset \cdots \subset R_k$ be nonempty open rectangles in $E_n$ with sides parallel to the axes and centered at 0. Suppose $S \subset E_n$ is a bounded set and, for each $x$ in S, let $i(x)$ be an integer satisfying $1 \leq i(x) \leq k$. If
\[
R^{x}=\{u\in E_{n}\colon u\;-\;x\in R_{i(x)}\}
\]

then there exist a finite number of points, $x_1, x_2, \ldots, x_l$, of $S$ such that $S \subset R^{x_1} \cup \cdots \cup R^{x_l}$, and each $v$ in $E_n$ belongs to at most $2^n$ of the sets $R^{x_1}, \ldots, R^{x_l}$.
\end{lemma}

\textsc{Proof.}\ Choose $x_1$ so that $i(x_1)$ is the largest possible (we can do this since $i(x) \leq k$). Having done so, we then choose $x_2 \in S - R^{x_1}$ so that $i(x_2)$ is largest possible and continue in this way by choosing $x_3 \in S - \{R^{x_1} \cup R^{x_2}\}$ so that $i(x_3)$ is maximal, etc. We thus obtain a sequence of open rectangles $R^{x_1}, R^{x_2}, \ldots$ centered at $x_1$, $x_2, \ldots$ with the property that each center $x_j$ belongs to no other rectangle $R^{x_i}, i \neq j$. But this means that the distance between $x_i$ and $x_j$, $i \neq j$, must exceed half the length of the

shortest side of $R_1$ (the smallest rectangle). Since, by assumption, this half-length is positive and the set $S$ is bounded it follows that the sequence $x_1, x_2, \ldots$ is finite and, after reaching the, say, $l$-th stage of our process we have obtained a covering $\{R^{x_1}, \ldots, R^{x_l}\}$ of $S$.

It now remains to be shown that a point $v \in E_n$ can belong to no more than $2^n$ of the rectangles $R^{x_1}, \ldots, R^{x_t}$. In order to see this, consider the hyperplanes through $v$ that are parallel to the coordinate hyperplanes. These form $2^n$ ``octants'' with vertex $v$ that cover all of $E_n$. Suppose two of the points, $x_i$ and $x_j$ ( $1 \leq i, j \leq l$), belong to the same octant in such a way that $R^{x_i}$ and $R^{x_t}$ both contain $v$. Since each of these two rectangles is a translate of one of the rectangles $R_1, \ldots, R_k$, one of these two, say $R^{x_i}$, must have sides that are at least as long as those of the other, $R^{x_i}$. But this, together with the fact that $v \in R^{x_i} \cap R^{x_t}$, implies that $x_j \in R^{x_i}$, contradicting the fact that each center $x_j$ belongs to no other rectangle $R^{x_i}$. Thus, at most one rectangle with center in any given ``octant'' can contain $v$ and the lemma is proved.

\begin{theorem}
\label{thm:c2-3-4}
There exists a constant $c = c(n)$, depending only on the dimension, such that whenever $f \in L^{1}(E_{n})$ and
\[
F_{s}=\{x\in E_{n}\colon m_{f}(x)>s>0\}
\]

then
\[
|F_{s}|\leq\frac{c\|f\|_{1}}{s},
\]

where $|F_s|$ denotes the Lebesgue measure of $F_s$. In particular, $m_f(x) < \infty$ for almost every $x \in E_n$.
\end{theorem}

\textsc{Proof.}\ We shall show that

\begin{equation}\label{eq:c2-3-5}
\left|\{x\in E_{n}\colon\overline{{m}}_{f}(x)\;>\;s\}\right|\;\leq\;2^{n}\frac{\|f\|_{1}}{s}
\end{equation}

for each $s > 0$. Since, by \eqref{eq:c2-3-2}, $F_s \subset \{x \in E_n: \overline{m}_f(x) > s/A_n\}$, we would then have

\[
\left|F_{s}\right|\leq\left|\left\{x\in E_{n}\colon\overline{m}_{f}(x)>\frac{s}{A_{n}}\right\}\right|\leq\frac{2^{n}A_{n}\|f\|_{1}}{s}
\]

and the theorem follows with $c = 2^{n}A_{n}$.

Let $S$ be any compact subset of $\{x \in E_n: \overline{m}_j(x) > s\}$. For each $x \in S$, therefore, there exists a cube $Q_x$ such that

\[
\frac{1}{\left|Q_{x}\right|}\int\limits_{Q_{x}}^{*}\left|f(t)\right|dt>s.
\]

By continuity of the integral, there exists a neighborhood $U$ of $x$ such that for each $u \in U$

\[
\frac{1}{\left|Q_{u}\right|}\int\limits_{Q_{u}}\left|f(t)\right|dt>s,
\]

where $Q_u = \{(w - x) + u: w \in Q_x\}$ is the cube of the same size as $Q_x$ that is centered at u. Since $S$ is compact there are a finite number of such neighborhoods covering $S$. Hence, there are a finite number of cubes $Q_1 \subset \cdots \subset Q_k$, centered at 0, such that for each $x \in S$ there exists an integer $i(x)$, $1 \leq i(x) \leq k$, for which

\[
\frac{1}{\left|Q^{x}\right|}\int\limits_{Q^{x}}\left|f(t)\right|dt>s,
\]

where $Q^x = \{u \in E_n: u - x \in Q_{i(x)}\}$. By Lemma \ref{lem:c2-3-3}, we can find a finite number of points $x_1, \ldots, x_l$ in $S$ such that

\[
s=Q^{x_{1}}\cup\cdots\cup Q^{x_{1}}
\]

and each $v\in E_n$ belongs to at most $2^n$ of these cubes. If, for a set $A$, $\chi_A$ denotes the characteristic function of $A$, this last property is equivalent to the inequality

\[
\sum_{j=1}^{l}\chi_{Q^{x_{j}}}\leq2^{n}\chi_{\bigcup_{j=1}^{l}Q^{x_{j}}}.
\]

Thus,

\[
\begin{aligned}|S|\leq&\left|\bigcup_{j=1}^{l}Q^{x_{j}}\right|\leq\sum_{j=1}^{l}|Q^{x_{j}}|\leq\frac{1}{s}\sum_{j=1}^{l}\int_{Q^{x_{j}}}|f(t)|dt\\ &=\frac{1}{s}\int_{E_{n}}\left\{\sum_{j=1}^{l}\chi_{_{Q^{x_{j}}}}(t)\right\}|f(t)|dt\\&\leq\frac{2^{n}}{s}\int_{E_{n}}\chi_{_{\bigcup_{j=1}^{l}Q^{x_{j}}}(t)}|f(t)|dt\\&\leq\frac{2^{n}}{s}\left\|f\right\|_{1}.\end{aligned}
\]

Since $S$ is an arbitrary compact subset of $\{x \in E_n: \overline{m}_f(x) > s\}$ inequality \eqref{eq:c2-3-5} and, hence, the theorem now follow immediately.

We have already observed that $m_g(x) \leq \|g\|_\infty < \infty$ for any $g \in L^\infty(E_n)$. From Theorem \ref{thm:c2-3-4} we have $m_f(x) < \infty$ for almost every $x$ when $f \in L^1(E_n)$. It follows from these two facts that $m_h(x) < \infty$ a.e. whenever $h \in L^p(E_n)$, $1 < p < \infty$. In order to see this we first observe that the maximal operator is an example of a sublinear\index{sublinear operator@sublinear operator} operator. By this we mean an operator $T$ mapping a linear space of measurable functions defined on a measure space into measurable functions defined on (possibly) another measure space satisfying

\begin{equation}
\tag{i}
|[T(f+g)](x)|\leq|(T f)(x)|+|(T g)(x)|
\end{equation}

almost everywhere (this property is called subadditivity\index{subadditivity@subadditivity}), and

\begin{equation}
\tag{ii}
|T(af)|=|a||Tf|
\end{equation}

whenever $a$ is a scalar and $f$ is in the domain of $T$. Thus, if, for $h \in L^{p}(E_{n})$, we put

\[
g(x)=\left\{\begin{aligned}{}&{{}h(x)}&{\quad\text{when }|h(x)| \leq1}\\ {}&{{}0}&{\quad\text{when }|h(x)| >1}\\ \end{aligned}\right.
\]

and $f = h - g$, then $g \in L^\infty(E_n)$ and $f \in L^1(E_n)$. By property (i) (subadditivity), therefore,

\[
m_{h}(x)=m_{f+g}(x)\leq m_{f}(x)+m_{g}(x).
\]

Consequently, $m_h(x) < \infty$ for almost every x. All these observations are, of course, equally valid for the cubic maximal function.

The maximal operator, as a mapping defined on $L^1(E_n)$ is not a bounded transformation into $L^1(E_n)$ (e.g., if $f$ is the characteristic function of the interval $[0,1]$ then $m_f$ is not even integrable). Theorem \ref{thm:c2-3-4}, however, tells us that it is ``almost'' bounded. In order to make this statement more precise we introduce the notion of the distribution function\index{distribution function@distribution function} $\lambda = \lambda_g$ of a function $g \in L^p(E_n)$: if $s > 0$ then $\lambda(s)$ is the Lebesgue measure of the set $F_s = \{x \in E_n: |g(x)| > s\}$.\footnote{The distribution function of g can obviously be defined more generally for any measurable function on an arbitrary measure space. In Chapter V we study many of the properties of these more general distribution functions.} We thus obtain a nonincreasing function $\lambda$ on the positive real axis. It is clear that $s^{p}\lambda(s) = \int_{F_s} s^p \, dx \leq \int_{F_s} |g|^p \, dx \leq \|g\|_p$. That is, $\lambda(s)$ is majorized by a constant times $s^{-p}$. This condition is clearly not sufficient to insure the integrability of $|g|^p$. On the other hand, the equality

\begin{equation}\label{eq:c2-3-6}
\Vert g\Vert_{p}^{p}\;=\;-\int\limits_{0}^{\infty}s^{p}\;d\lambda(s)\;=\;p\;\int\limits_{0}^{\infty}s^{p-1}\lambda(s)\;d s,
\end{equation}

which is immediate for $g$ simple (the general case then follows by approximating $|g|$ from below by simple functions), indicates that the condition $\lambda(s) \leq$ (constant) $s^{-p}$ is close to being equivalent with the finiteness of $\int |g|^{p}$. Theorem \ref{thm:c2-3-4} tells us that $m_f$ has a distribution function $\lambda$ satisfying $\lambda(s) \leq c \|f\|_1 s^{-1}$. It is in this sense that we say that the maximal operator is almost bounded as an operator on $L^1(E_n)$.

The maximal operator, however, is bounded when restricted to $L^p(E_n)$. $1 < p \leq \infty$:

\begin{theorem}
\label{thm:c2-3-7}
There exists a constant $b = b(p, n)$, depending on the dimension $n$ and the exponent $p > 1$, such that
\[
\left\|m_{f}\right\|_{p}\leq b\|f\|_{p}
\]

for all $f \in L^{p}(E_{n})$.
\end{theorem}

\textsc{Proof.}\ Suppose $f \in L^p(E_n)$, $1 < p \leq \infty$, and $s > 0$. For each $x \in E_n$, let

\[
f^{s}(x)=\left\{\begin{aligned}&f(x)&\text{ if }|f(x)|>s\\ &0&\text{ if }|f(x)|\leq s\end{aligned}\right.
\]

and

\[
f_{s}(x)=\left\{\begin{matrix}{0}&{\quad\text{if }|f(x)|>s}\\ {f(x)}&{\quad\text{if }|f(x)|\leq s.}\\ \end{matrix}\right.
\]

Then $f = f^s + f_s$, with $f^s \in L^1(E_n)$ and $f_s \in L^\infty(E_n)$.

We let $\lambda$, $\lambda^{s}$, $\lambda_{s}$ be the distribution functions of $m_{f}$, $m_{f^{s}}$, and $m_{f_{s}}$. Since, by subadditivity, $m_{f} \leq m_{f^{s}} + m_{f_{s}}$ it follows that $\lambda(2s) \leq \lambda^{s}(s) + \lambda_{s}(s)$. We already observed that the maximal function of an essentially bounded function $g$ is bounded by $\|g\|_{\infty}$. Thus, $m_{f_{s}}(x) \leq \|f_{s}\|_{\infty} \leq s$ and it follows that $\lambda_{s}(s) = 0$. Therefore $\lambda(2s) \leq \lambda^{s}(s)$.

From \eqref{eq:c2-3-6}, this last inequality and Theorem \ref{thm:c2-3-4} we therefore obtain

\[
\begin{aligned}\|m_{f}\|_{p}^{p}&=p2^{p}\int\limits_{0}^{\infty}s^{p-1}\lambda(2s)ds\leq p2^{p}\int\limits_{0}^{\infty}s^{p-1}\lambda^{s}(s)ds\\&\leq p2^{p}\int\limits_{0}^{\infty}s^{p-1}cs^{-1}\Biggl\{\int\limits_{E_{n}}|f^{s}(x)|dx\Biggr\}ds\\&=p2^{p}c\int\limits_{0}^{\infty}s^{p-2}\Biggl\{\int\limits_{|f(x)|>s}|f(x)|dx\Biggr\}ds\\&=p2^{p}c\int\limits_{E_{n}}|f(x)|\Biggl\{\int\limits_{0}^{|f(x)|}s^{p-2}ds\Biggr\}dx\\&=\frac{p2^{p}c}{p-1}\int\limits_{E_{n}}|f(x)||f(x)|^{p-1}dx,\end{aligned}
\]

which proves the theorem with $b = b(p, n) = 2[pc/(p - 1)]^{1/p}$.

Another way of expressing $m_{f}$ is the following one:

Let $\varphi$ be the characteristic function of the solid unit sphere $\{x \in E_n: |x| \leq 1\}$ divided by $\Omega_n$ (so that $\int_{E_n} \varphi(x) \, dx = 1$). Let $\varepsilon > 0$ and put $\varphi_\varepsilon(t) =$

$\varepsilon^{-n}\varphi(t/\varepsilon)$. Then, if $f \in L^p(E_n)$, $1 \leq p \leq \infty$, and if, for the sake of notational simplicity, we assume $f \geq 0$, we have

\[
(f*\varphi_{\varepsilon})(x)=\frac{1}{\varepsilon^{n}\Omega_{n}}\int\limits_{|t|\leq\varepsilon}f(x-t)\;d t;
\]

therefore,

\begin{equation}\label{eq:c2-3-8}
m_{f}(x)=\sup_{\varepsilon>0}(f*\varphi_{\varepsilon})(x).
\end{equation}

Similarly, if $\varphi$ is the characteristic function of the cube $\{x = (x_1, \ldots, x_n) \in E_n: -\frac{1}{2} \leq x_j \leq \frac{1}{2}, j = 1, 2, \ldots, n\}$, then the right side of \eqref{eq:c2-3-8} gives us $\overline{m}_f(x)$. We have seen that $\overline{m}_f$ and $m_f$ are ``equivalent'' in the sense that inequalities \eqref{eq:c2-3-2} and \eqref{eq:c2-3-2p} hold. It is natural, therefore, to ask about the relation of $m_f$ with the supremum in \eqref{eq:c2-3-8} when $\varphi \geq 0$ is a more general function.

For example, suppose $\varphi = \sum_{k=1}^{m} c_k \chi_k$, $c_k > 0$, where $\chi_k$ is the characteristic function of the sphere $\sigma_k = \{x \in E_n: |x| \leq r_k\}$. Then, for $f \geq 0$,

\[
\begin{aligned}(f*\varphi_\varepsilon)(x)&=\sum_{k=1}^{m} c_{k}\varepsilon^{-n}\int\limits_{|t|\leq\varepsilon r_{k}}f(x-t)dt\\&=\sum_{k=1}^{m} c_{k} r_{k}^{n}\Omega_{n}\frac{1}{\Omega_{n}(\varepsilon r_{k})^{n}}\int\limits_{|t|\leq\varepsilon r_{k}}f(x-t)dt\\&\leq m_{f}(x)\sum_{k=1}^{m} c_{k}|\sigma_{k}|=m_{f}(x)\|\varphi\|_{1}.\end{aligned}
\]

That is.

\begin{equation}\label{eq:c2-3-9}
\sup_{\varepsilon>0}(f*\varphi_{\varepsilon})(x)\;\leq\;m_{f}(x)\|\varphi\|_{1}.
\end{equation}

This result clearly extends to any $\varphi \in L^1(E_n)$ of the form $\varphi(t) = \psi(|t|)$ (i.e., $\varphi$ is radial) with $\psi$ a nonnegative decreasing function on $[0, \infty)$. We need only approximate $\varphi$ from below by an ascending sequence of simple functions of the type just considered and apply the Lebesgue monotone convergence theorem.

In particular, if we choose $\varphi(t) = c_n/(|t|^2 + 1)^{(n+1)/2}$ then $u(x, y) = (f * \varphi_y)(x), y > 0$, is the Poisson integral of f and we have the theorem:

\begin{theorem}
\label{thm:c2-3-10}
If $u(x, y), y > 0$, is the Poisson integral of $f \in L^{p}(E_{n})$, $1 \leq p \leq \infty$, then $|u(x, y)| \leq m_{f}(x)$.
\end{theorem}

It is not hard to establish the converse to this last result:

\begin{theorem}
\label{thm:c2-3-11}
If $f \geq 0$ and its Poisson integral
\[
u(x,y)=\int\limits_{E_{n}}f(x-t)P(t,y)dt,\qquad y>0,
\]

is defined (for example, if $f$ belongs to $L^p(E_n)$, $1 \leq p \leq \infty$), then there exists a constant $A = A(n)$, depending only on the dimension, such that
\[
m_{f}(x)\leq A\left\{\sup_{y>0}u(x,y)\right\}.
\]

\end{theorem}

\textsc{Proof.}\ Fix r > 0, then

\[
\begin{aligned}\sup_{y>0}u(x,y)\geq u(x,r)&=c_{n}\int\limits_{E_{n}}f(x-t)\frac{r}{(|t|^{2}+r^{2})^{(n+1)/2}}dt\\&\geq c_{n}^{\prime}\int\limits_{|t|\leq r}f(x-t)r^{-n}dt\\&=c_{n}^{\prime}\Omega_{n}\Biggl\{\frac{1}{\Omega_{n} r^{n}}\int\limits_{|t|\leq r}f(x-t)dt\Biggr\}.\end{aligned}
\]

Thus, taking the supremum over all $r > 0$ of the expressions on the right we obtain the theorem with $A(n) = 1/c_n' \Omega_n$.

The importance of the maximal function operator lies in the fact that it majorizes many important operators in analysis. We have just seen an example of this situation in the case of the Poisson integral. The following theorem gives us a method for using this majorization property in order to obtain pointwise convergence results for many of the families of operators that we shall encounter.

\begin{theorem}
\label{thm:c2-3-12}
Let $\{T_\varepsilon\}$, $0 < \varepsilon$, be a family of linear operators mapping $L^p(E_n)$, $1 \leq p \leq \infty$, into the space of measurable functions on $E_\pi$. For each $h \in L^p(E_n)$ define $M h$ by letting $(M h)(x) = \sup_{\varepsilon > 0} |(T_\varepsilon h)(x)|$ for each $x \in E_n$.
\end{theorem}

Suppose that there exists a constant $a > 0$ and a real number $q \geq 1$ such that

\begin{equation}\label{eq:c2-3-13}
\left|\left\{x:\left(M h\right)(x)>\lambda\right\}\right|\leq(a\|h\|_{p}\lambda^{-1})^{q}
\end{equation}

for all $\lambda > 0$ and $h \in L^p(E_n)$. If there exists a dense subset $\mathcal{D}$ of $L^p(E_n)$ such that $\lim_{\varepsilon \to 0} (T_\varepsilon g)(x)$ exists and is finite almost everywhere whenever $g \in \mathcal{D}$, then for each $f \in L^p(E_n)$, $\lim_{\varepsilon \to 0} (T_\varepsilon f)(x)$ exists and is finite almost everywhere.

\textsc{Proof.}\ Let $f$ be a function in $L^p(E_n)$. For each positive $k$, let $F_k$ be the set of all $x \in E_n$ such that $|\langle T_{\varepsilon}, f \rangle(x) - (T_{\varepsilon}, f)(x) \rangle| > 2k^{-1}$ for infinitely many pairs $(\varepsilon', \varepsilon')$ tending to $(0, 0)$. For each $\eta > 0$ we can find $g$ and $h$ in $L^p(E_n)$ such that $f = g + h$ with $g \in \mathcal{D}$ and $\|h\|_p < \eta$. We shall show that $|F_k| \leq (2a(k + 1)\eta)^q$; since $\eta$ is arbitrary this means that $F_k$ has measure 0. Thus, $F = \bigcup_{k=1}^{\infty} F_k$ has measure 0 and, for $x \in E_n - F$, $\lim_{\varepsilon \to 0} (T_{\varepsilon} f)(x)$ exists and is finite. This, then, would prove the theorem.

Let $G$ be the set of all $x \in E_n$ for which $\lim_{\varepsilon \to 0} (T_\varepsilon g)(x)$ exists and is finite. Since $E_n - G$ has measure zero it suffices to show that $|F_k \cap G| \leq (2a(k + 1)\eta)^q$. Since $(T_\varepsilon f)(x) - (T_\varepsilon f)(x) = \{(T_\varepsilon, h)(x) - (T_\varepsilon^c h)(x)\} + \{(T_\varepsilon, g)(x) - (T_\varepsilon^c g)(x)\}$ and, for $x \in G$, $\lim_{\varepsilon', \varepsilon'' \to 0} \{(T_\varepsilon g)(x) - (T_\varepsilon^c g)(x)\} = 0$, it follows that $F_k \cap G$ is contained in the set of all $x \in E_n$ such that $|(T_\varepsilon h)(x) - (T_\varepsilon^c h)(x)| \geq k^{-1}$ for infinitely many pairs $(\varepsilon', \varepsilon')$ tending to $(0, 0)$. For each such $x$ we must certainly have $(Mh)(x) \geq (2k)^{-1}$; therefore, $F_k \cap G \subset \{x \in E_n; (Mh)(x) > 2^{-1}(k + 1)^{-1}\}$. The desired inequality $|F_k \cap G| \leq (a\eta^2(k + 1))^q$ now follows from our hypothesis.

An example of such a family is given by defining $T_{\varepsilon}f$, for $\varepsilon>0$, by letting

\[
(T_{\varepsilon}f)(x)=(\varepsilon^{n}\Omega_{n})^{-1}\int\limits_{|t|<\varepsilon}f(x-t)d t
\]

for each $x \in E_n$. In this case we have $Mf = m_f$ and, by Theorem \ref{thm:c2-3-4}, we have

\[
\left|\left\{x:\left(M f\right)(x)>\lambda\right\}\right|\leq c\left\|f\right\|_{1}\lambda^{-1}
\]

for $f\in L^{1}(E_{n})$. Furthermore, if $g$ belongs to the class of continuous functions with compact support (a dense subset of $L^{1}(E_{n}))$ then it is clear that $\lim_{\varepsilon\to0}(T_{\varepsilon}g)(x)=g(x)$. Thus, by Theorem \ref{thm:c2-3-12}, $\lim_{\varepsilon\to0}(T_{\varepsilon}f)(x)$ exists and is finite for almost all $x$. From this we easily deduce that almost every $x$ is a point of differentiability of the integral of $f$ (see Chapter I, \eqref{eq:c1-1-23} and \eqref{eq:c1-1-24}); more precisely, we have:

\begin{corollary}
\label{cor:c2-3-14}
Iff is locally integrable on $E_{n}$ (i.e., f is integrable over any bounded subset of $E_{n}$) then
\[
\lim_{\varepsilon\to0}\varepsilon^{-n}\int\limits_{|t|<\varepsilon}\left[f(x-t)-f(x)\right]d t=0
\]

for almost every $x \in E_n$.
\end{corollary}

\textsc{Proof.}\ Since the problem is local we can assume $f \in L^1(E_n)$ (otherwise we can multiply f by the characteristic function of a sphere of radius $r$ about 0 and obtain the desired convergence at almost every point within this sphere; letting $r$ tend to $\infty$ we then obtain it for almost all points of $E_n$). We already know that

\[
\lim_{\varepsilon\to0}\left(T_{\varepsilon}f\right)(x)=\lim_{\varepsilon\to0}\left(\varepsilon^{n}\Omega_{n}\right)^{-1}\int\limits_{|t|<\varepsilon}f(x-t)dt
\]

exists and is finite for almost all x. We only need to show that this limit\index{non-tangential limit@non-tangential limit} is $f(x)$ a.e. But

\[
\begin{aligned}\left\|T_{\varepsilon}f-f\right\|_{1}&=\int\limits_{E_{n}}\left|\frac{1}{\varepsilon^{n}\Omega_{n}}\int\limits_{|t|<\varepsilon}\left[f(x-t)-f(x)\right]dt\right|dx\\&\leq\frac{1}{\varepsilon^{n}\Omega_{n}}\int\limits_{|t|<\varepsilon}\left\{\int\limits_{E_{n}}\left|f(x-t)-f(x)\right|dx\right\}dt\to0\end{aligned}
\]

as $\varepsilon \to 0$ since the $L^1$ modulus of continuity $\omega_{1,f}(t) \to 0$ as $|t| \to 0$ (see the proof of Theorem \ref{thm:c1-1-18} in Chapter I). Consequently, there exists a sequence, $\{\varepsilon_k\}$, tending to 0 such that $(T_{\varepsilon_k}f)(x) \to f(x)$ a.e. as $k \to \infty$. It follows, therefore, that $\lim_{t \to 0}(T_{\varepsilon}f)(x) = f(x)$ a.e.

This method extends immediately to Poisson integrals and we can obtain from it a direct proof of Theorem \ref{thm:c2-2-1}. Our intention, however, is to extend this theorem to include the more general nontangential approach to the boundary that was mentioned briefly at the beginning of this section. We first make this notion more precise and then will show that Poisson integrals have these more general boundary values.

If $\alpha > 0$ and $x_0 \in E_n$ the cone\index{cone@cone!in En+1@in $E_{n+1}$} in $E_{n+1}^+$ with vertex $(x_0, 0)$ and aperture $\alpha$ is the region

\[
\Gamma_{\alpha}(x_{0})=\{(x,y)\in E_{n+1}^{+}\colon|x-x_{0}|<\alpha y\}.
\]

If $u$ is defined in $E_{n+1}^{+}$ we shall say that it has a nontangential limit $l$ at $x_0 \in E_n$ if, for each $\alpha > 0$,

\begin{equation}\label{eq:c2-3-15}
\lim\limits_{(x,y)\to(x_{0},0)}u(x,y)=l,
\end{equation}

where $(x,y)$ tends to $(x_{0},0)$ within the cone $\Gamma_{\alpha}(x_{0})$.

The importance of the nontangential approach to the boundary will become clearer later. In several basic results that we shall develop, the vertical approach to the boundary is not sufficient (see, e.g. \eqref{frp:c2-5-4}).

\begin{theorem}
\label{thm:c2-3-16}
If $f \in L^p(E_n)$, $1 \leq p \leq \infty$, then the Poisson integral
\[
\int\limits_{E_{n}}f(x-t)P(t,y)dt=u(x,y)
\]

has the nontangential limit $f(x_0)$ at each $x_0$ belonging to the Lebesgue set of $f$; in particular, it has these limits at almost all points of $E_n$.
\end{theorem}

\textsc{Proof.}\ We first observe that the proof of Theorem \ref{thm:c1-1-25} in the first chapter consisted in showing that

\[
\begin{aligned}0&=\lim_{y\to0}\int\limits_{E_{n}}\left|f(x_{0}-t)-f(x_{0})\right|\left|\varphi_{y}(t)\right|dt\\&=\lim_{y\to0}\int\limits_{E_{n}}\left|f(t)-f(x_{0})\right|\left|\varphi_{y}(x_{0}-t)\right|dt\end{aligned}
\]

whenever $x_{0}$ belongs to the Lebesgue set of $f$. In case $\varphi_{y}(t)$ is the Poisson kernel $P(t,y)=c_{n}[y/(y^{2}+|t|^{2})^{(n+1)/2}]$ we have the easily established inequality

\begin{equation}\label{eq:c2-3-17}
P(x-t,y)=\varphi_{y}(x-t)\leq d_{\alpha}\varphi_{y}(x_{0}-t)=d_{\alpha}P(x_{0}-t,y)
\end{equation}

for $(x, y) \in \Gamma_\alpha(x_0)$, where $d_\alpha^{2/(n+1)} = \max\{1 + 2\alpha^2, 2\}$. Consequently, using Lemma \ref{lem:c1-1-17}, part (b) (in Chapter I),

\[
\begin{aligned}|u(x,y)-f(x_{0})|&=\left|\int\limits_{E_{n}}f(t)P(x-t,y)dt-f(x_{0})\int\limits_{E_{n}}P(x-t,y)dt\right|\\&\leq\int\limits_{E_{n}}|f(t)-f(x_{0})|P(x-t,y)dt\\&\leq d_{\alpha}\int\limits_{E_{n}}|f(t)-f(x_{0})|P(x_{0}-t,y)dt\\&=d_{\alpha}\int\limits_{E_{n}}|f(t)-f(x_{0})|\varphi_{y}(x_{0}-t)dt.\end{aligned}
\]

But, since the last term tends to 0 as y tends to 0, the theorem is proved.

An immediate consequence of the result we have just proved is that the Poisson integral $u$ is nontangentially bounded\index{non-tangentially bounded@non-tangentially bounded} at almost every point of $E_n$. That is, it is a function $u$ defined in $E_{n+1}^+$ that is bounded in the truncated cones $\Gamma_a(x_0) \cap \{(x, y) \in E_{n+1}; 0 < y \leq 1\}$ for almost all $x_0 \in E_n$ (we truncate these cones at height 1; since the functions we shall consider are continuous in $E_{n+1}^+$ the height at which we truncate these cones is of no importance). In fact, the maximal function $m_f$ can be used to give the following quantitative expression of this:

\begin{equation}\label{eq:c2-3-18}
\sup_{(x,y)\in\Gamma_{\alpha}(x_{0})}\left|u(x,y)\right|\leq d_{\alpha}m_{f}(x_{0}).
\end{equation}

This is a simple consequence of Theorem \ref{thm:c2-3-10} and inequality 3.17 involving the Poisson kernel.

The following theorem (a basic tool used in the theory of the $H^{p}$ spaces developed in Chapter VI) shows that, for harmonic functions, nontangential limits and nontangential boundedness are essentially equivalent.

\begin{theorem}
\label{thm:c2-3-19}
Suppose u is harmonic in $E_{n+1}^+$ and is nontangentially bounded at all the points of a subset $S \subset E_n$ of positive measure, then u has nontangential limits at almost every point of S.
\end{theorem}

\textsc{Proof.}\ First, we observe that it suffices to show that whenever $S \subset E_n$ is a subset of positive measure at which $u$ is nontangentially bounded then there exists a subset of $S$, also of positive measure, at whose points $u$ has nontangential limits. Considering only cones $\Gamma_\alpha(q)$ with rational apertures $\alpha$ we see that $S$ is a denumerable union of subsets $F = F_{\alpha,m}$ such that $|u(x, y)| \leq m$ for all

\[
(x,y)\in\left\{\bigcup_{q\in F}\Gamma_{\alpha}(q)\right\}\cap\{p=(x,y);0<y<2\}=\mathcal{A},
\]

$m = 1, 2, \ldots\footnote{Throughout this proof we shall consistently use this notation: the points of $E_{n}$ (i.e., those points of $E_{n+1}$ with last coordinate 0) will be denoted by the letter $q$. The letter $p$ will be used to denote points of $E_{n+1}^{+}$.}$ That is, $u$ is uniformly bounded in the region $\mathcal{A}$ formed by the union of congruent vertical cones whose vertices are the points of $F$ and are bounded above by the hyperplane $y = 2$.

Now, it is easy to see that it is sufficient for us to show that $u(p)$ tends to a limit when $p$ tends to $q$ within $\Gamma_a(q)$, for almost all $q \in F$. For, if we discard the sets where this does not occur, for each $F = F_{\alpha,m}$ (recall that $\alpha$ is rational and $m$ is a positive integer), we will have discarded a set of measure zero; if $q$ is a point of $S$ not in this set of measure 0 and $\beta > 0$ any rational aperture, then there exists a positive integer $k$ such that $|u(p)| = |u(x, y)| \leq k$ for $p \in \Gamma_\beta(q)$, $0 < y < 2$; that is, $q \in F_{\beta,k}$. Consequently, $u(p)$ tends to a limit as $p$ tends to $q$ within $\Gamma_\beta(q)$. Thus, a non-tangential limit exists for $u$ at all such $q's$.

A further reduction is obtained by observing that, dividing by $m$, we can assume that $|u| \leq 1$ in $\mathcal{A}$. Moreover, we can also assume that $S$ is contained in a cube of side 1; since it suffices to prove the theorem for each intersection of $S$ with a member of a mesh of cubes of side 1.

Having made all these simplifications, we let $\mathcal{D} = \{\bigcup_{q \in F} \Gamma_a(q)\} \cap \{p = (x, y); 0 < y < 1\}, \mathcal{B} = \partial \mathcal{D}$ (= boundary of $\mathcal{D}$), $\mathcal{D}_j$ denote the set $\mathcal{A}$ translated by -1/j in the direction of the $y$ axis and $\mathcal{G}_j = \mathcal{D}_j \cap E_n$. Moreover, we put $u_j(x, y) = u(x, y + 1/j)$, the characteristic function of $\mathcal{G}_j$ will be denoted by $\chi_j$, the Poisson integral of $\chi_j(x)u_j(x, y)$ by $\varphi_j(x, y)$ and $\psi_j(x, y) = u_j(x, y) - \varphi_j(x, y)$. We observe that $|\varphi_j| \leq 1$ (this follows from Lemma \ref{lem:c1-1-17}, part (b), of Chapter I and the fact that the Poisson kernel is nonnegative).

The sequence of functions $f_j(s) = \chi_j(s)u_j(s, 0)$ is uniformly bounded in the $L^2(E_n)$ norm (for $\|f_j\|_2 \leq \|\chi_j\|_2 = (\int_{\mathcal{G}_j} ds)^{1/2} \leq$ (volume of cube of side $1 + 2\alpha/j)^{1/2} \leq (1 + 2\alpha)^{n/2}$). Thus, there exists a subsequence $\{f_{j_k}\}$ that

converges weakly to a function $f \in L^2(E_n)$. In particular,

\[
\varphi_{f_{k}}(x,y)=\int\limits_{E_{n}}f_{f_{k}}(s)P(x-s,y)ds\to\int\limits_{E_{n}}f(s)P(x-s,y)ds=\varphi(x,y)
\]

as $k \to \infty$, for each $(x, y) \in E_{n+1}^+$.

Since it is obviously true that $\lim_{j\to\infty}u_{j}(x,y)=u(x,y)$, we have the existence of the limit

\[
\psi(x,y)=\lim_{k\to\infty}\psi_{j_{k}}(x,y)=u(x,y)-\varphi(x,y).
\]

Since $\varphi$ is the Poisson integral of a function in $L^{2}(E_{n})$ it has, by Theorem \ref{thm:c2-3-16}, nontangential limits at almost all points of $E_{n}$. Since $u = \varphi + \psi$ we shall, therefore, want to show that $\lim \psi(p) = 0$ when $p$ tends to $q$ within $\Gamma_{\alpha}(q)$, for almost all $q \in F$.

Toward this end we first observe that the functions $\psi$ satisfy

(1) $|\psi_j(p)| \leq 2$ for $p \in \mathcal{D}$ (since $|\psi_j| \leq |u_j| + |\varphi_j| \leq 1 + 1 = 2$ in $\mathcal{D}$);

(2) $\psi_j(p) \to 0$ as $p \to q \in F$, $p \in \mathcal{D}$ (in fact, it follows from Theorem \ref{thm:c2-3-16} that this occurs at all interior points of $\mathcal{G}_j$).

Now we claim that our theorem would be established if we could find a function w, harmonic in $E_{n+1}^{+}$, satisfying

(a) $w(x, y) \geq 0$ in $E_{n+1}^+$;

(b) $w(x, y) \geq 2 \text{ in } \mathcal{B} - F;$

(c) w has nontangential limits 0 at almost every point of F.

If this were the case, $w(p) \pm \psi_j(p) \geq 0$ in $\mathcal{B} - F$ because of (1) and (b). Furthermore, (2) and (a) imply

\[
\lim_{p\to q\in{\cal F},p\in\mathcal{D}}\{w(p)\pm\psi_{j}(p)\}\geq\lim_{p\to q\in{\cal F},p\in\mathcal{D}}w(p)\geq0.
\]

It follows that $w(p) \pm \psi_j(p) \geq 0$ for all $p \in \mathcal{D}$. For, if not, the minimum principle implies the existence of a sequence $\{p_k\} \subset \mathcal{D}$, tending to a point of $F$ as $k \to \infty$, such that $w(p_k) \pm \psi_j(p_k) < -\varepsilon$ for each $k$, where $\varepsilon$ is some positive number. But, by (2), $\psi_j(p_k) \to 0$ as $k \to \infty$ and, thus, for $k$ large $w(p_k)$ would be negative, contradicting property (a). Consequently, letting $j$ tend to $\infty$ (through a subsequence), we obtain $w(p) \pm \psi(p) \geq 0$ for all $p \in \mathcal{D}$. That is, $|\psi(p)| \leq w(p)$ for all $p \in \mathcal{D}$. Hence, the desired result, $\lim \psi(p) = 0$ when $p$ tends to $q$ within $\Gamma_a(q)$, for almost all $q \in F$, is a consequence of property (c).

We shall now continue the proof by constructing a function w satisfying

the three properties (a), (b), and (c). Let $\xi$ be the characteristic function of $E_{n} - F$ and let, for y > 0,

\[
w(x,y)=2y+c y\int_{E_{n}}\frac{\xi(s)}{(y^{2}+|x-s|^{2})^{(n+1)/2}}d s,
\]

where $c$ will be a positive constant (to be chosen later). It is clear that $w \geq 0$ in $E_{n+1}^{+}$. Property (c) is an immediate consequence of Theorem \ref{thm:c2-3-16}. Thus, we need only show that $w$ satisfies property (b). We first observe that if $(x, y) \in \mathcal{B} - F$ with $y = 1$ then $w(x, y) \geq 2y = 2$. Now suppose that $p = (x, y) \in \mathcal{B}$ with $0 < y < 1$. Consider the inverted cone with vertex $p$ and aperture $\alpha$ (see figure). If $C$ is the interior of the sphere

\begin{figure}[htbp]
\centering
\figcone
\caption{The cone with vertex $p$ and aperture $\alpha$}
\label{fig:c2-p78}
\end{figure}

formed by the intersection of this cone with the hyperplane $E_n$ then $C$ and $F$ are disjoint. For, if this were not the case, there would exist a point $q \in C \cap F$ and, consequently, $p$ would be in the interior of $\Gamma_\alpha(q)$. But this is impossible since $p \in \mathcal{B}$. Thus,

\[
\begin{aligned}w(p)&\geq cy\int\limits_{C}\frac{ds}{(y^{2}+|x-s|^{2})^{(n+1)/2}}=c\omega_{n-1}y\int\limits_{0}^{\alpha y}\frac{r^{n-1}dr}{(y^{2}+r^{2})^{(n+1)/2}}\\&=c\omega_{n-1}\int\limits_{0}^{\alpha}\frac{r^{n-1}dr}{(1+r^{2})^{(n+1)/2}}.\end{aligned}
\]

This last expression can obviously be made to exceed 2 by choosing c appropriately.

Having constructed the function $w$ we have, therefore, shown that $\psi(p)$ has the limit $0$ as $p$ tends to $q$ within $\Gamma_a(q)$, for almost every $q$ in $F$. We next claim that the majorization $|\psi| \leq w$ in $\bigcup_{q\in F} \Gamma_a(q)$ implies that for almost every $q$ in $F$, $|\psi(p)| \leq w(p)$, where $p$ may lie in a cone with vertex $q$ of arbitrary aperture, as long as $p$ is sufficiently close to $q$. A simple geometric argument shows that this is, indeed, the case for each

$q$ which is a point of density\index{point of density@point of density} of $F$.\footnote{Let $\xi_F = \xi$ be the characteristic function of a measurable set $F \subset E_n$, then $x$ is a point of density of $F$ if $\Omega_n^{-1} r^{-n} \int_{|x-t| < r} \xi(t) \, dt \to 1$ as $r \to 0$. By Corollary \ref{cor:c2-3-14} almost every point of $F$ is a point of density.} Consequently, $\lim \psi(p) = 0$ as $p$ tends to $q \in F$ nontangentially for almost every $q$ in $F$. This means that $u$ and $\varphi$ have the same nontangential limits at almost every $q$ in $F$ and, therefore, the theorem is proved.

We shall be interested in extending this result to functions that are harmonic in sets of variables (these will be called multiply harmonic\index{multiply harmonic functions@multiply harmonic functions} functions). By this we mean twice differentiable functions, $u$, defined on a domain, $\mathcal{D}$, of points $(x^{(1)}, x^{(2)}, \ldots, x^{(k)}) = (x_1^{(1)}, \ldots, x_{n_1}^{(1)}, x_1^{(2)}, \ldots, x_{n_2}^{(2)}, \ldots, x_1^{(k)}, \ldots, x_{n_k}^{(k)})$ belonging to the Cartesian product $E_{n_1} \times E_{n_2} \times \cdots \times E_{n_k}$ of $k$ Euclidean spaces, that are harmonic in each of the variables $x^{(t)}, 1 \leq j \leq k$; that is,

\[
\sum_{i=1}^{n_{f}}\frac{\partial^{2}u}{(\partial x_{i}^{(j)})^{2}}=0
\]

for $1 \leq j \leq k$. We shall encounter such functions in the next chapter, where we deal with analytic functions of several variables. In this case, each of the spaces $E_{n_j}$ will be two-dimensional and can be identified with a copy of the complex plane.

It is not hard to state this extension and to indicate the changes in the proof of Theorem \ref{thm:c2-3-19} that are needed in order to establish it. One of these changes involves results analogous to Theorems \ref{thm:c2-2-5} and 3.16 for ``iterated Poisson integrals'' (i.e., solutions of a variant of the Dirichlet problem involving multiply harmonic functions).

These iterated Poisson integrals are examples of a family of operators somewhat more general than that occurring in Theorem \ref{thm:c2-3-12}, in that they involve several independent parameters tending to 0. That is, the family $\{T_e\}$ is parametrized by the set $\{\varepsilon = (\varepsilon_1, \varepsilon_2, \ldots, \varepsilon_m) \in E_m \colon \varepsilon_j > 0, j = 1, 2, \ldots, m\}$ (the first octant\index{octant@octant} in $E_m$); their behavior as $\varepsilon$ tends to $0 \in E_m$ will be the object of our study. In the case of iterated Poisson integrals let us, for simplicity consider the case when $m = 2$. That is, suppose we identify $E_n$ with the Cartesian product $E_{n_1}$ and $E_{n_2}$, where $n_1 + n_2 = n$, and let

\[
P_{j}(t,\varepsilon_{j})=c_{n_{j}}\frac{\varepsilon_{j}}{(|t|^{2}+\varepsilon_{j}^{2})^{(n_{f}+1)/2}}
\]

for $t = (t_1, \ldots, t_{n_j}) \in E_{n_j}$, $j = 1, 2$. Then, for $f \in L^p(E_n)$ we define $T_\varepsilon f = T_{(\varepsilon_1, \varepsilon_2)} f$ by letting

\[
(T_{\varepsilon}f)(x)=\int_{E_{n_{1}}}P_{1}(x^{(1)}-t,\varepsilon_{1})\Biggl\{\int_{E_{n_{2}}}P_{2}(x^{(2)}-s,\varepsilon_{2})f(t,s)d s\Biggr\}d t,
\]

where

\[
x=\left(x_{1}^{(1)},\ \ldots,\ x_{n_{1}}^{(1)},\ x_{1}^{(2)},\ \ldots,\ x_{n_{2}}^{(2)}\right)=\left(x^{(1)},x^{(2)}\right)\in E_{n_{1}}\times E_{n_{2}}=E_{n}.
\]

If $f(x) = f_1(x^{(1)})f_2(x^{(2)})$ with $f_j \in L^p(E_{n_j}) \cap C_0(E_{n_j})$, $j = 1, 2$, it follows from Theorem \ref{thm:c2-2-1} that $(T_{\varepsilon}f)(x)$ converges to $f(x)$ uniformly as $\varepsilon = (\varepsilon_1, \varepsilon_2) \to (0, 0)$. The class of finite linear combinations of such functions is dense in $L^p(E_n)$, $1 \leq p < \infty$; thus, if we can show that inequality \eqref{eq:c2-3-13} holds we can then conclude that $\lim_{\varepsilon \to 0} (T_{\varepsilon}f)(x)$ exists and is finite almost everywhere. It is not hard to establish the following stronger result for $1 < p \leq \infty$ for $(Mf)(x) = \sup_{\varepsilon_1, \varepsilon_2 > 0} |(T_{(\varepsilon_1, \varepsilon_2)} f)(x)|$

\begin{equation}\label{eq:c2-3-20}
\left\|M f\right\|_{p}\leq A\|f\|_{p},
\end{equation}

where $A$ depends on $p$, $n_1$ and $n_2$, but is independent of $f \in L^p(E_n)$ (see the discussion preceding \eqref{eq:c2-3-6}). In order to do this we define $M^{(1)}$ to be the operator of $L^p(E_n)$ assigning to $f$ the maximal function of $f(\cdot, x^{(2)})$; that is,

\[
(M^{(1)}f)(x)=(M^{(1)}f)(x^{(1)},x^{(2)})=\sup_{r>0}\Omega_{n_{1}}^{-1}r^{-n_{1}}\int\limits_{|t|\leq r}\left|f(x^{(1)}-t,x^{(2)})\right|d t.
\]

Similarly, we let

\[
(M^{(2)}f)(x)=\sup_{r>0}\Omega_{n_{2}}^{-1}r^{-n_{2}}\int\limits_{|s|\leq r}\left|f(x^{(1)},x^{(2)}-s)\right|d s.
\]

Thus, by Theorem \ref{thm:c2-3-7}.

\begin{equation}\label{eq:c2-3-21}
\begin{aligned}{\|M^{(2)}M^{(1)}f\|_{p}}&{{}\leq b(p,n_{2})\|M^{(1)}f\|_{p}\leq b(p,n_{2})b(p,n_{1})\|f\|_{p}}\\ {}&{{}=B\|f\|_{p}.}\\ \end{aligned}
\end{equation}

On the other hand, it follows from \eqref{thm:c2-3-10} that

\[
\begin{aligned}|(T_{\varepsilon}f)(x)|&=\left|(T_{(\varepsilon_{1},\varepsilon_{2})}f)(x)\right|\\&\leq\int\limits_{E_{\mathbb{R}^{2}}}P_{2}(x^{(2)}-s,\varepsilon_{2})(M^{(1)}f)(x^{(1)},s)\,ds\\&\leq(M^{(2)}M^{(1)}f)(x^{(1)},x^{(2)}).\end{aligned}
\]

Consequently, $(Mf)(x) = \sup_{\varepsilon_1, \varepsilon_2 > 0} |(T_{(\varepsilon_1, \varepsilon_2)} f)(x)| \leq (M^{(2)} M^{(1)} f)(x)$ and inequality \eqref{eq:c2-3-20} follows from \eqref{eq:c2-3-21}.

We can, therefore, conclude that $\lim (T_\varepsilon f)(x)$ as $\varepsilon = (\varepsilon_1, \varepsilon_2) \to (0, 0)$ exists and is finite for almost all $x \in E_n$. It is clear that we can modify the above argument to apply to $m$-fold iterated Poisson integrals. More precisely, we have the following theorem:

\begin{theorem}
\label{thm:c2-3-22}
If $f \in L^p(E_n)$, $1 < p < \infty$, and
\[
\begin{aligned}&u(x,\varepsilon)=u(x^{(1)},x^{(2)},\ldots,x^{(m)};\varepsilon_{1},\varepsilon_{2},\ldots,\varepsilon_{m})=\\ &\int\limits_{E_{n_{1}}}\cdots\int\limits_{E_{n_{m}}}P_{1}(x^{(1)}-t^{(1)},\varepsilon_{1})\cdots P_{m}(x^{(m)}-t^{(m)},\varepsilon_{m})\\&\qquad f(t^{(1)},\ldots,t^{(m)})dt^{(1)}\cdots dt^{(m)},\end{aligned}
\]

where $t^{(j)}$ and $x^{(j)}$ belong to $E_{n_j}$ and $\varepsilon = (\varepsilon_1, \varepsilon_2, \ldots, \varepsilon_m)$ belong to the first octant of $E_m$ (that is, $\varepsilon_j > 0$ for $j = 1,2, \ldots, m$), then
\[
\lim_{\varepsilon\to0}u(x,\varepsilon)=\lim_{\varepsilon_{1},\ldots,\varepsilon_{m}\to0}u(x^{(1)},x^{(2)},\ldots,x^{(m)};\varepsilon_{1},\ldots,\varepsilon_{m})=f(x)
\]

for almost every $x \in E_n$.
\end{theorem}

\textsc{Proof.}\ We have already seen that $\lim_{\varepsilon \to 0} u(x, \varepsilon)$ exists and is finite almost everywhere. We need only show that this limit is $f(x)$ a.e. As in the proof of Corollary \ref{cor:c2-3-14} it suffices to show that $\|u(\cdot, \varepsilon) - f\|_{p} \to 0$ as $\varepsilon \to 0$. This can be done in much the same way as was done in the proof of Theorem \ref{thm:c1-1-18} in Chapter I. In order to simplify the notation we shall, again, present the proof when m = 2; the method extends in an obvious way to the general case. We have, by Lemma \ref{lem:c1-1-17}, part (b), of the first chapter and by a change of variables:

\[
\begin{aligned}u(x,\varepsilon)&-f(x)\\&=\int\limits_{E_{n_{1}}E_{n_{2}}}\int P_{1}(t,\varepsilon_{1})P_{2}(s,\varepsilon_{2})\\&\qquad\{f(x^{(1)}-t,x^{(2)}-s)-f(x^{(1)},x^{(2)})\}dt ds\\&=\int\limits_{E_{n_{1}}}\int\limits_{E_{n_{2}}}P_{1}(t,1)P_{2}(s,1)\\&\qquad\{f(x^{(1)}-\varepsilon_{1}t,x^{(2)}-\varepsilon_{2}s)-f(x^{(1)},x^{(2)})\}dt ds.\end{aligned}
\]

\[
\begin{aligned}\text{For }(h,k)&\in E_{n_{1}}\times E_{n_{2}}\text{ let}\\\omega_{p,f}(h,k)&=\omega(h,k)\\&=\left(\int\limits_{E_{n_{1}}}\int\limits_{E_{n_{2}}}\right.\\&\qquad\left.|f(x^{(1)}+h,x^{(2)}+k)-f(x^{(1)},x^{(2)})|^{p} dx^{(1)}dx^{(2)}\right)^{1/p}\end{aligned}
\]

be the $L^p$ modulus of continuity of $f$. We have seen that $\omega(h, k) \to 0$ as $|h|, |k| \to 0$ and that $\omega(h, k) \leq 2 \| f\|_p$. Thus, by Minkowski's inequality for integrals and the Lebesgue dominated convergence theorem,

\[
\begin{aligned}&\left(\int\limits_{E_{n}}|u(x,\varepsilon)-f(x)|^{p} dx\right)^{1/p}\\&=\int\limits_{E_{n_{1}}}\int\limits_{E_{n_{2}}}\omega(-\varepsilon_{1}t,-\varepsilon_{2}s)P_{1}(t,1)P_{2}(s,1)dt ds\rightarrow0\end{aligned}
\]

as $\varepsilon_{1}, \varepsilon_{2} \rightarrow 0$. This completes the proof.

A modification of the argument we have just given shows that the theorem (as well as its nontangential version that follows below) is valid also for $p = \infty$. Since, in any case, this is a consequence of Theorem \ref{thm:c2-3-24} proved below, we need not present this modification.

Theorem \ref{thm:c2-3-22} asserts that the iterated Poisson integral of a function in $L^{p}(E_{n})$ converges to the function almost everywhere if we approach the points $x_{0} = (x_{0}^{(1)}, x_{0}^{(2)}, \ldots, x_{0}^{(m)})$ in $E_{n}$ ``downward'' along the products of the lines $x^{(j)} = x_{0}^{(j)}$ in $E_{n_j+1}$, $j = 1, 2, \ldots, m$ (compare with the discussion at the beginning of this section). As was the case before, more general limits exist. These are generalizations of nontangential limits and are defined in the following way: Suppose $n_1 + n_2 + \cdots + n_k = n$ and

\[
\Gamma_{\alpha_{j}}^{(j)}(x_{0}^{(j)})=\{(x^{(j)},y_{j})\in E_{n_{j}+1}^{+};|x^{(j)}-x_{0}^{(j)}|<\alpha_{j}y_{j}\},\qquad j=1,2,\ldots,k.
\]

A function, $u$, defined on $E_{n_1+1}^+ \times \cdots \times E_{n_k+1}^+$, is said to have the nontangential limit $l$ in each set of variables\index{non-tangential limit@non-tangential limit!in a set of variables@in a set of variables}\footnote{We are grouping the $n$ variables $x_1, x_2, \ldots, x_n$ in $k$ mutually disjoint sets. The $j$-th set consists of the variables $x_{N_j-1+1}, \ldots, x_{N_j}$, where $N_i = n_1 + n_2 + \cdots + n_l$ for $1 \leq l \leq k$. The concept we have just defined involves this grouping and we should, properly, refer to it in the definition. On the other hand, we shall see that in all our applications it will be clear which grouping is involved; therefore, we avoid referring to it for the sake of simplicity.} at the point $x_0 = (x_0^{(1)}, x_0^{(2)}, \ldots, x_0^{(k)})$ if $u(x^{(1)}, y_1; x^{(2)}, y_2; \ldots; x^{(k)}, y_k) = u(x^{(1)}, x^{(2)}, \ldots, x^{(k)}; y_1, y_2, \ldots, y_k)$ tends to $l$ as the point $(x^{(1)}, y_1; x^{(2)}, y_2; \ldots; x^{(k)}, y_k)$ tends to $x_0$ within the Cartesian product $\Gamma_{\alpha_1}^{(1)} \times \Gamma_{\alpha_2}^{(2)} \times \cdots \times \Gamma_{\alpha_k}^{(k)}$ for all $k$-tuples ($\alpha_1, \alpha_2, \ldots, \alpha_k$) of positive real numbers (i.e., each of the points $(x^{(j)}, y_j)$ tends nontangentially to $x_0^{(j)}, 1 \leq j \leq k$). By the argument used to prove Theorem \ref{thm:c2-3-22} and the inequality \eqref{eq:c2-3-17} applied to each of the kernels $P_{j_1}, 1 \leq j \leq m$, we then obtain\index{non-tangentially bounded@non-tangentially bounded!in a set of variables@in a set of variables}

\begin{corollary}
\label{cor:c2-3-23}
If $f \in L^p(E_n)$, $1 < p < \infty$, and $u(x, \epsilon)$ is the function defined in Theorem \eqref{thm:c2-3-22}, then $u(x, \epsilon)$ has the nontangential limit $f(x_0)$ in each set of variables at almost all $x_0 \in E_n$.
\end{corollary}

We say that $u$, defined on $E_{n_1+1} \times E_{n_2+1} \times \cdots \times E_{n_{k+1}}$ is nontangentially bounded in each set of variables at $x_0 \in E_n$ if $u$ is bounded in the Cartesian product of the truncated cones $\Gamma_{\alpha_j}(x_0^{(j)}) \cap \{(x^{(j)}, y_j) \in E_{n_j+1}; 0 < y_j \leq 1\}$, $1 \leq j \leq k$, for all $k$-tuples $(\alpha_1, \ldots, \alpha_k)$ of positive real numbers. Suppose $u$ is the iterated Poisson integral of a function $f \in L^p(E_n)$, $1 < p < \infty$. The fact that, at almost all points of $E_n$, $u$ is nontangentially bounded in each set of variables is a somewhat more elementary result than Corollary \ref{cor:c2-3-23} (it is, of course, an immediate consequence of that corollary).

In complete analogy to the case of harmonic functions on $E_{n+1}^{+}$, this nontangential boundedness is almost everywhere equivalent to the existence of nontangential limits in each set of variables, even for harmonic

functions that are not Poisson integrals. More precisely, we have the following extension of Theorem \ref{thm:c2-3-19}:

\begin{theorem}
\label{thm:c2-3-24}
Suppose $n_1 + n_2 + \cdots + n_k = n$ and $u$ is a function defined in the Cartesian product $E_{n_1+1}^+ \times E_{n_2+1}^+ \times \cdots \times E_{n_k+1}^+$ that is harmonic in each set of variables $(x^{(j)}, y_j) = (x_1^{(j)}, \ldots, x_{n_j}^{(j)}, y_j)$ in $E_{n_j+1}^+$, $1 \leq j \leq k$. (That is, $u$ is twice differentiable and satisfies
\[
\Delta_{j}u=\sum_{i=1}^{n_{j}}\frac{\partial^{2}u}{(\partial x_{i}^{(j)})^{2}}+\frac{\partial^{2}u}{\partial y_{j}^{2}}=0).
\]

\end{theorem}

If $u$ is nontangentially bounded in each set of variables at all the points of a subset $S \subset E_n$ of positive measure, then $u$ has nontangential limits in each set of variables at almost every point of $S$.

\textsc{Proof.}\ The main ideas in establishing this theorem were developed in the proof of Theorem \ref{thm:c2-3-19}. Therefore, we shall only indicate the changes in that proof that are needed to obtain the present theorem. We confine ourselves to the case k=2 and $n_{1}=n_{2}=n$; (the argument extends easily to the general case).

In analogy with the notation used in the proof of Theorem \ref{thm:c2-3-19} we let $p = (p_1, p_2) = (x^{(1)}, y_1; x^{(2)}, y_2)$ denote the general point of $E_{n+1}^{+} \times E_{n+1}^{+}$; and let $q = (q_1, q_2) = (x^{(1)}, 0; x^{(2)}, 0)$ denote the general point of $E_n \times E_n = E_{2n}$, (the latter is the distinguished boundary\index{distinguished boundary@distinguished boundary} of $E_{n+1}^{+} \times E_{n+1}^{+}$).\footnote{The terminology ``distinguished boundary'' is often used in the literature; it indicates that we are considering only a part of the topological boundary of $E_{n+1}^{+}\times E_{n+1}^{+}$.} By reasoning similar to that used before, we essentially reduce the problem to showing that if $|u| \leq 1$ in the region

\[
\mathcal{A}=\left\{\bigcup_{(q_{1},q_{2})\in F}\gamma_{\alpha}(q_{1},q_{2})\right\}\cap\{p=(x^{(1)},y_{1};x^{(2)},y_{2});0<y_{1},y_{2}<2\},
\]

then $u(p)$ tends to a limit when $p$ tends to $q = (q_1, q_2)$ within $\gamma_\alpha(q) = \gamma_\alpha(q_1, q_2)$, for almost all of $q \in F \subset E_{2n}$. Here we have written $\gamma_\alpha(q_1, q_2)$ for the Cartesian product cone $\Gamma_{\alpha_1}(q_1) \times \Gamma_{\alpha_2}(q_2)$. By taking appropriate subsets we can assume that $F$ is closed and, in addition, that $F$ is contained in the unit sphere of $E_{2n}$.

We let

\[
\mathcal{D}=\left\{\bigcup_{q\in F}\gamma_{\alpha}(q)\right\}\cap\{p=(x^{(1)},y_{1};x^{(2)},y_{2});0<y_{1},y_{2}<1\},
\]

$\mathcal{A}$, denote the set translated by -1/j in the direction of the $y_1$ and $y_2$ axes, and $\mathcal{G}_j = \mathcal{A}_j \cap E_{2n}$. Moreover, we put $u_j(p) = u(x^{(1)}, y_1 + 1/j)$;

$x^{(2)}, y_2 + 1/j)$; the characteristic function of $\mathcal{G}_j$ will be denoted by $\chi_j$, the iterated Poisson integral of $\chi_j(x^{(1)}; x^{(2)})u_j(x^{(1)}, 0; x^{(2)}, 0)$ by $\varphi_j$ and $\psi_j = u_j - \varphi_j$.

As was the case in the proof of Theorem \ref{thm:c2-3-19} we can find a sequence $\{\varphi_{jk}\}$ converging to an iterated Poisson integral (of a function in $L^{2}(E_{2n}))$, at each point of $E_{n+1}^{+}\times E_{n+1}^{+}$. We denote this limit by $\varphi$. Since $u_j\to u$ as $j\to\infty$, it follows that $\psi_{jk}(p)=u_{jk}(p)-\varphi_{jk}(p)$ tends to the limit $\psi(p)=u(p)-\varphi(p)$, for each $p\in E_{n+1}^{+}\times E_{n+1}^{+}$. Since $\varphi$ is the Poisson integral of a function in $L^{2}(E_{2n})$, by Corollary \ref{cor:c2-3-23} it has nontangential limits in each set of variables at almost all points of $E_{2n}$. We shall show that $\lim\psi(p)=0$, when $p$ tends to $q$ within $\gamma_{\alpha}(q)$, for almost all $q\in F$.

In keeping with the argument used in the proof of Theorem \ref{thm:c2-3-19} the crucial point consists in constructing a multiply harmonic function $w \in E_{n+1}^{+} \times E_{n+1}^{+}$, satisfying

(a) $w \geq 0$ in $E_{n+1}^{+} \times E_{n+1}^{+}$,

(b) $\liminf_{p' \to p} w(p') \geq |\psi_i(p)|, j = 1, 2, \ldots$ wherever $p$ is a point of the boundary of $\mathcal{D}$,

(c) w has zero nontangential limits in each set of variables at almost every point of F.

Let $\xi$ be the characteristic function of that part of $E_{2n} - F$ contained in a fixed sphere of sufficiently large radius centered at the origin (the lower bounds required for the size of this sphere are more easily described in the context of the proof and will be obvious). We then define $w = 2(y_1 + y_2) + cW(x^{(1)}, y_1; x^{(2)}, y_2)$, where $W$ is the repeated Poisson integral of $\xi$, and $c$ is a positive constant to be determined later.

The condition (a) is obvious, and condition (c) follows from Corollary \ref{cor:c2-3-23} (since $\xi$ belongs to all the spaces $L^p(E_{2n})$, $1 < p < \infty$). Thus, we need only prove (b). But $\psi_j(p) = 0$, for $p \in F$, so inequality (b) is true for $p \in F$. We consider separately the following remaining parts of the boundary of $\mathcal{D}$:

(i) points where either $y_{1}$ or $y_{2}$ is 1;

(ii) points where $0 < y_{1} < 1$, and $0 < y_{2} < 1$;

(iii) points such that $y_{1} > 0$ and $y_{2} = 0$, or such that $y_{2} > 0$ and $y_{1} = 0$.

Since $|\psi_i| \leq 2$, property (b) is obvious in case (i) holds.

Suppose $p = (p_1, p_2) = (x^{(1)}, y_1; x^{(2)}, y_2)$ satisfies (ii). Then consider the Cartesian product of the inverted cones with vertices $p_1$ and $p_2$ and apertures $\alpha_1$ and $\alpha_2$. These cones intersect the hyperplanes $y_1 = 0$, and $y_2 = 0$ (in each $E_{n+1}^+$) in two open spheres $C_1$ and $C_2$. We claim that

$C_1 \times C_2$ does not contain a point of $F$. (If this were not the case $p$ would belong to $\mathcal{D}$.) Thus $\xi$ equals 1 at all points $C_1 \times C_2$, provided we chose the sphere containing $F$ which was used to define $\xi$ to be large enough to include these Cartesian products. The spheres, $C_1$ and $C_2$ are centered at $x^{(1)}$ and $x^{(2)}$ and have radii $\alpha_1 y_1$ and $\alpha_2 y_2$, respectively. Therefore,

\[
\begin{aligned}w(p)&\geq cW(p)\geq c\prod_{j=1}^{2}\left(c_{n} y_{j}\int\limits_{C_{j}}\frac{ds^{(j)}}{(y_{j}^{2}+|x^{(j)}-s^{(j)}|^{2})^{(n+1)/2}}\right)\\&=c\prod_{j=1}^{2}\left(y_{j} c_{n}\omega_{n-1}\int\limits_{0}^{\alpha_{j} y_{j}}\frac{r^{n-1}dr}{(y_{j}^{2}+r^{2})^{(n+1)/2}}\right)\geq2,\end{aligned}
\]

if c is large enough.

Now suppose $p = (p_1, p_2) = (x^{(1)}, y_1; x^{(2)}, 0)$ satisfies the condition (iii), (the case $y_1 = 0$, $y_2 > 0$ is similar) in this case $w(p)$ is not defined and so we must estimate $\lim \inf w(p')$, as $p' \in \mathcal{D}$ approaches $p$. Letting $p' = (p_1', p_2')$, and writing $y_1'$ and $y_2'$ for the $y$ coordinates of $p_1', p_2$, respectively we have

\[
\begin{aligned}&w(p^{\prime})\geq2y_{1}^{\prime}+c c_{n}^{2}\intop_{E_{n}}\frac{y_{1}^{\prime}}{\left|p_{1}^{\prime}-q_{1}\right|^{(n+1)}}\\&\qquad\left\{\intop_{E_{n}}\frac{y_{2}^{\prime}}{\left|p_{2}^{\prime}-q_{2}\right|^{(n+1)/2}}\xi(q_{1},q_{2})d q_{2}\right\}d q_{1}.\end{aligned}
\]

It follows from Fatou's lemma that

$\liminf_{p^{\prime}\to p,p^{\prime}\in\mathcal{D}}w(p^{\prime})\geq$

\[
\begin{aligned}&2y_{1}\:+\:c c_{n}^{2}\int_{E_{n}}\frac{y_{1}}{\left|{p_{1}\:-\:q_{1}}\right|^{(n+1)}}\overline{{\lim\inf_{p_{2}^{^{\prime}}\to p_{2}}c_{n}}}\\&\qquad\int_{E_{n}}\frac{y_{2}^{^{\prime}}}{\left|{p_{2}^{^{\prime}}\:-\:q_{2}}\right|^{(n+1)/2}}\;\xi(q_{1}q_{2})\;d q_{2}\Bigg\}\;d q_{1}.\end{aligned}
\]

We now claim that

\begin{equation}\label{eq:c2-3-25}
\lim_{p_{2}^{\prime}\to p_{2}}\inf c_{n}\int_{E_{n}}\frac{y_{2}^{\prime}}{\left|p_{2}^{\prime}-q_{2}\right|^{(n+1)}}\xi(q_{1},q_{2})d q_{2}\geq\xi(q_{1},p_{2}).
\end{equation}

To see this we first note that the inequality certainly holds when $\xi(q_1, p_2) = 0$, since the integrand on the left is nonnegative. The only other value assumed by $\xi$ is 1 on the open set which is the intersection of $E_{2n} - F$ with an open sphere in $E_{2n}$. Thus $\xi$ is continuous at these points. It follows that the integral \eqref{eq:c2-3-25} actually converges to $\xi(q, p_2)$ (this can be shown by the same argument that was used to establish the last part of (b) in Theorem \ref{thm:c2-2-1}). Thus \eqref{eq:c2-3-25} is proved and, consequently,

\[
\lim_{p^{\prime}\to p,p^{\prime}\in\mathcal{D}}\displaystyle w(p^{\prime})\geq2y_{1}+c c_{n}\int_{E_{n}}\frac{y_{1}}{\left|{p_{1}-q_{1}}\right|^{(n+1)/2}}\;\xi(q_{1},p_{2})\;d q_{1}.
\]

The rest of the proof is then essentially a repetition of the last part of the argument we used in order to establish Theorem \ref{thm:c2-3-19} applied to the functions $\psi_{j}(\bar{p}_{1}, p_{2})$ and

\[
v(\bar{p}_{1})=2y_{1}+cc_{n}\int\limits_{E_{n}}\frac{y_{1}}{|\bar{p}_{1}-q_{1}|^{(n+1)}}\xi(q_{1},p_{2})dq_{1}
\]

in the region $\mathcal{D}_{p_2}$ consisting of those $\bar{p}_1 \in E_{n+1}^+$ that are contained in some truncated cone $\Gamma_{\alpha_1}(q_1)$ with $(q_1, p_2) \in F$. (At this stage $p_2$ is a fixed point which lies on the boundary $E_n$ of $E_{n+1}^+$.) This completes the proof of (b), and with it the majorization

\begin{equation}\label{eq:c2-3-26}
|\psi(p)|\;\leq\;w(p),p\in\mathcal{D}.
\end{equation}

To conclude the proof of the theorem, we shall need to extend \eqref{eq:c2-3-26} so that for almost all $q \in F$, $|\psi(p)| \leq w(p)$ where $p$ may lie in product cones $\Gamma_{\beta_1}(q_1) \times \Gamma_{\beta_2}(q_2)$, with arbitrary apertures $\beta_1$ and $\beta_2$, as long as $p$ is sufficiently close to $q = (q_1, q_2)$. This will indeed be the case at all points of strong density\index{point of density@point of density!strong@strong} of $F$. These are the points which satisfy the conclusion of the following lemma:\index{strong density@strong density!for points in En@for points in $E_n$}

\begin{lemma}
\label{lem:c2-3-27}
Let $\xi$ be the characteristic function of an arbitrary measurable set $F$ of $E_{2n} = E_n \times E_n$. Then for almost every $x = (x^{(1)}, x^{(2)}) \in F$
\[
\lim_{h_{1}\to0\atop h_{2}\to0}\Omega_{n}^{-2}h_{1}^{-n}h_{2}^{-n}\int\limits_{|s^{(1)}|<h_{1},|s^{(2)}|<h_{2}}\int\limits_{\xi(x^{(1)}-s^{(1)},x^{(2)}-s^{(2)})}d s^{(1)}d s^{(2)}=1.
\]

\end{lemma}

\textsc{Proof.}\ For any $f \in L^{p}(E_{n} \times E_{n})$, we define the strong maximal function $\mu_{f}$, as

\[
\mu_{f}(x^{(1)},x^{(2)})=
\]

\[
\begin{aligned}&\sup_{h>0,\;h_{2}>0}\;\Omega_{n}^{-2}h_{1}^{-n}h_{2}^{-n}\int\int_{|s^{(1)}|<h_{1},|s^{(2)}|<h_{2}}\\&\qquad|f(x^{(1)}-s^{(1)},x^{(2)}-s^{(2)})|d s^{(1)}d s^{(2)}.\end{aligned}
\]

Then, as in the proof of \eqref{eq:c2-3-20} $\mu_f(x) \leq (M^{(2)} M^{(1)} f)(x)$, and so

\begin{equation}\label{eq:c2-3-28}
\|\mu_{j}(x)\|_{p}\;\leq\;A_{p}\|f\|_{p},\qquad1<p<\infty.
\end{equation}

By the same argument as in the proof of Theorem \ref{thm:c2-3-22} it then follows that

\begin{equation}\label{eq:c2-3-29}
\begin{aligned}\lim_{h_{1}\to0\atop h_{2}\to0}\Omega_{n}^{-2}h_{1}^{-n}h_{2}^{-n}\iint\limits_{|s^{(1)}|<h_{1},|s^{(2)}|<h_{2}}f(x^{(1)}-s^{(1)},x^{(2)}-s^{(2)}) ds^{(1)} ds^{(2)}\\=f(x^{(1)},x^{(2)})\end{aligned}
\end{equation}

for almost every $(x^{(1)}, x^{(2)}) \in E_{2n}$, for any fixed $f \in L^p(E_{2n})$, $1 < p < \infty$. The validity of the assertion \eqref{eq:c2-3-29}, however, depends only on the values of $f$ in a neighborhood of the point $(x^{(1)}, x^{(2)})$, and so \eqref{eq:c2-3-29} holds almost

everywhere when $f$ is merely locally in $L^p(E_{2n})$, $1 < p < \infty$. Thus it holds for bounded functions and, in particular, for the function $\xi$. This proves Lemma \ref{lem:c2-3-27}.

We can now conclude the proof of Theorem \ref{thm:c2-3-24}. Suppose that $(q_1, q_2)$ is a point of strong density of $F$. For fixed $y_1, y_2 > 0$, consider the set

\[
A=\{(x^{(1)},x^{(2)})\colon|x^{(1)}-q_{1}|<\beta_{1}y_{1},|x^{(2)}-q_{2}|<\beta_{2}y_{2}\}.
\]

For any $(x^{(1)}, x^{(2)}) \in A$ we also consider the set

\[
B=\{\left(\bar{x}^{(1)},\bar{x}^{(2)}\right);\left|x^{(1)}-\bar{x}^{(1)}\right|<\alpha_{1}y_{1},\left|x^{(2)}-\bar{x}^{(2)}\right|<\alpha_{2}y_{2}\}.
\]

If $y_1$ and $y_2$ are sufficiently small, then $B$ must intersect $F$, otherwise we would obtain a contradiction with the fact that $q = (q_1, q_2)$ is a point of strong density of $F$. Thus, if $p = (p_1, p_2)$, $p_1 = (x^{(1)}, y_1)$, $p_2 = (x^{(2)}, y_2)$, and $p$ is sufficiently close to $q$, then $p \in \Gamma_{\beta_1}(q_1) \times \Gamma_{\beta_2}(q_2)$ implies that $p \in \Gamma_{\alpha_1}(\bar{x}^{(1)}) \times \Gamma_{\alpha_2}(\bar{x}^{(2)})$, with $(\bar{x}^{(1)}, \bar{x}^{(2)}) \in F$.

The estimate \eqref{eq:c2-3-26} then also holds for all $p \in \Gamma_{\beta_1}(q_1) \times \Gamma_{\beta_2}(q_2)$, if $p$ is sufficiently close to $q = (q_1, q_2)$ and the desired nontangential limit of $\psi$ is therefore established.

\section{Subharmonic Functions and Majorization by Harmonic Functions}

In later chapters we will need to bound from above, with appropriate functions that are harmonic, functions having the form $|F|^p$, where $F$ is either a holomorphic function of several complex variables (see Chapter III) or a vector-valued function whose components are harmonic (see Chapter VI). These functions $|F|^p$ will turn out to be subharmonic. We can give an indication of the nature of subharmonic functions and of this type of majorization by considering the banal case of functions of a single real variable. In this case the harmonic functions are the linear ones and, as will be apparent, the subharmonic functions are the convex functions. The latter are generally defined to be continuous \footnote{It can be shown that it suffices, in this one dimensional case, to suppose the function to be measurable. The general theory of subharmonic functions can be carried out for upper semicontinuous functions; however, though this is important in some applications we will have occasion to apply it only to continuous functions and, therefore, we shall limit ourselves to this somewhat (technically) simpler case.} functions $\varphi$, defined on an interval of the real line, satisfying

\begin{equation}\label{eq:c2-4-1}
\varphi\left(\frac{x+y}{2}\right)\leq\frac{\varphi(x)+\varphi(y)}{2}
\end{equation}

whenever the interval $[x, y]$ is contained in the domain of $\varphi$. This condition can easily be shown to be equivalent to the following one: If the interval $[x, y]$ is contained in the domain of $\varphi$ then the linear (harmonic)

function having the same values as $\varphi$ at the end points $x$ and $y$ majorizes\footnote{We say that a function $f$ majorizes a function $g$ if the two functions have the same domain $D$, are real valued, and $f(x) \geq g(x)$ for all $x \in D$.} $\varphi$ in this interval. In case the convex function $\varphi$ has a continuous second derivative at all points of its domain $D$, then \eqref{eq:c2-4-1} is equivalent to the condition $\varphi''(x) \geq 0$ for all $x \in D$.

In the first part of this section we generalize these facts and deal with functions defined on bounded domains; in the second part (as well as in later chapters) we consider these problems on certain unbounded domains that arise naturally in our study.

We say that a function, $s$, defined and continuous in a domain $\mathcal{D} \subset E_n$ is subharmonic\index{subharmonic function@subharmonic function} if

\begin{equation}\label{eq:c2-4-2}
s(x_{0})\leq\mathcal{M}_{x_{0},s}(r)=\frac{1}{\omega_{n-1}}\intop_{\Sigma_{n-1}}s(x_{0}+r t^{\prime})d t^{\prime}
\end{equation}

whenever the sphere $\{x \in E_n; |x - x_0| \leq r\}$ is contained in $\mathcal{D}$. This inequality clearly generalizes \eqref{eq:c2-4-1}.

The fact that a nonnegative second derivative implies convexity has the following extension to n dimensions:

\begin{theorem}
\label{thm:c2-4-3}
Suppose $s$ has continuous second partial derivatives in a domain $\mathcal{D}$ and $(\Delta s)(x) \geq 0$ for all $x \in \mathcal{D}$, then $s$ satisfies the mean value inequality \eqref{eq:c2-4-2} at all points of $\mathcal{D}$.
\end{theorem}

\textsc{Proof.}\ Suppose the closed solid sphere of radius $r_0$ about $x_0$ belongs to $\mathcal{D}$ and $\Sigma_\varepsilon, \Sigma_r, 0 < e < r \leq r_0$, denote the surfaces of the spheres of radii $\varepsilon$ and $r$ about $x_0$. Applying Green's theorem to the functions $s$ and

\[
v(x)=\left\{\begin{aligned}{}&{{}\left|x-x_{0}\right|^{-(n-2)}-r^{-(n-2)},\quad}&{}&{{}\text{if }n>2}\\ {}&{{}\left.-\log|x-x_{0}|\right.+\log r,\quad}&{}&{{}\text{if }n=2,}\\ \end{aligned}\right.
\]

in the shell domain, $S$, between $\Sigma_e$ and $\Sigma_r$ we obtain (since $\Delta v = 0$ and $\Delta s \geq 0$ in $S$)

\[
0\geq\int_{s}\int\left(s\Delta v-v\Delta s\right)d x=\int_{\partial s}\left(s\frac{\partial v}{\partial n}-v\frac{\partial s}{\partial n}\right)d\sigma,
\]

where $\partial S$ is the boundary of $S$, $\partial/\partial n$ denotes differentiation in the direction of the outward directed normal to $\partial S$, and $d\sigma$ is the element of surface area on $\partial S$. Thus, as in the proof of Theorem \ref{thm:c2-1-1},

\[
0\geq\left(\int_{\Sigma_{r}}+\int_{\Sigma_{\varepsilon}}\right)\frac{\partial v}{\partial n}s\;d\sigma-\left(\int_{\Sigma_{r}}+\int_{\Sigma_{\varepsilon}}\right)v\frac{\partial s}{\partial n}d\sigma.
\]

But $v$ is 0 on $\Sigma_r$ and $\int_{\Sigma_\varepsilon} \partial s / \partial n \, v \, d\sigma \leq 0$. (Here $\partial / \partial n$ denotes differentiation in the direction of the normal to $\Sigma_\varepsilon$ directed toward the center of this sphere. Thus, applying Green's theorem to the functions $s$ and 1 in the closed region $\{x; |x - x_0| \leq \varepsilon\}$ we see that $\int_{\Sigma_\varepsilon} \partial s / \partial n \, d\sigma \leq 0$; moreover, $v \geq 0$ on $\Sigma_\varepsilon$ whenever $\varepsilon \leq r$.) Hence, when $n > 2$,

\[
0\geq\left(\int_{\Sigma_{r}}-\int_{\Sigma_{\varepsilon}}\right)\left(-(n-2)|x-x_{0}|^{-(n-1)}s(x)\right)d\sigma(x)
\]

and, when n = 2,

\[
0\geq\left(\int_{\Sigma_{\tau}}-\int_{\Sigma_{\varepsilon}}\right)\left(-\left|x-x_{0}\right|^{-1}s(x)\right)d\sigma(x).
\]

In either case,

\[
\frac{1}{\omega_{n-1}\varepsilon^{n-1}}\;\int_{\Sigma_{\varepsilon}}s(x)\;d\sigma(x)\;\leq\;\frac{1}{\omega_{n-1}r^{n-1}}\int_{\Sigma_{r}}s(x)\;d\sigma(x)\;=\;\mathcal{M}_{x_{0},s}(r).
\]

Letting $\varepsilon$ tend to 0 the left side tends to $s(x_{0})$ and we obtain inequality \eqref{eq:c2-4-2}.

In later chapters we will want to show that certain functions are subharmonic. In general, they will not be twice differentiable; thus, Theorem \ref{thm:c2-4-3} cannot be applied in order to establish the desired subharmonicity. With the appropriate interpretation of the statement $\Delta s \geq 0$ (in the sense of distributions), Theorem \ref{thm:c2-4-3} remains valid for all continuous $s$. Further details can be found in \eqref{frp:c2-5-8} below. For our purposes, however, the following more specific variant of the theorem we just proved will be more useful.

\begin{theorem}
\label{thm:c2-4-4}
Suppose $s \geq 0$ is continuous in a domain $\mathcal{D}$, has continuous second partial derivatives in the subset $\mathcal{R} = \{x \in \mathcal{D}; s(x) > 0\}$ and satisfies $\Delta s \geq 0$ in $\mathcal{R}$. Then $s$ is a subharmonic function in $\mathcal{D}$.
\end{theorem}

\textsc{Proof.}\ We have to show that $s$ satisfies the mean value inequality \eqref{eq:c2-4-2}. Thus, let $x_0 \in \mathcal{D}$ and suppose the closed sphere $S_r = \{x; |x - x_0| \leq r\}$ is contained in $\mathcal{D}$. Let $u$ be the solution of the Dirichlet problem for the closed domain $S_r$ having boundary values $s(x)$ for $x \in \partial S_r = \{t; |t - x_0| = r\}$ (see Corollary \ref{cor:c2-1-11}). If we can show that $s(x) \leq u(x)$ for $x \in S_r$, inequality \eqref{eq:c2-4-2} is a consequence of the mean-value property (Theorem \ref{thm:c2-1-1}):

\[
\begin{aligned}s(x_{0})\leq u(x_{0})&=\mathcal{M}_{x_{0},u}(r)=\frac{1}{\omega_{n-1}}\int_{\Sigma}u(x_{0}+rt^{\prime})dt^{\prime}\\&=\frac{1}{\omega_{n-1}}\int_{\Sigma}s(x_{0}+rt^{\prime})dt^{\prime}.\end{aligned}
\]

Suppose, on the contrary, that $\sup_{x\in S_r}\{s(x)-u(x)\}=c>0$, and let $F=\{x\in S_r;\,s(x)-u(x)=c\}$. Then the set $F$ is obviously closed and is contained in the interior of $S_r$. On the other hand, $s(x)>0$ at each point of $F$ (since $u\geq0$). We can, therefore, apply Theorem \ref{thm:c2-4-3} to the function $s-u$ in the set $\mathcal{R}$ and we obtain, for each $x\in F(\subset\mathcal{R})$

\[
c=s(x)-u(x)\leq\frac{1}{\omega_{n-1}}\int\limits_{\Sigma}\left\{s(x+\rho t^{\prime})-u(x+\rho t^{\prime})\right\}d t^{\prime}
\]

for all $\rho > 0$ sufficiently small. Since $s - u \leq c$ in $S_r$ and $x$ is interior to $S_r$, the continuity of $s - u$ implies that $s(x + \rho t') - u(x + \rho t') = c$ for all $t' \in \Sigma$ and $\rho$ sufficiently small. Thus, $F$ is open. Since $S_r$ is connected $F$ must be either empty or all of $S_r$; but, as we have already observed, the boundary of $S_r$ is disjoint from $F$ and, consequently, $F$ is empty. This shows that $s(x) \leq u(x)$ for all $x$ in $S_r$ and the theorem is proved.

It is useful to note that we have actually proved the following more general result:

\begin{customthm}{4.4$'$}
\label{thm:c2-4-4p}
Suppose $s \geq 0$ is continuous in a domain $\mathcal{D}$ and is subharmonic in the set $\mathcal{R} = \{x \in \mathcal{D}; s(x) > 0\}$, then $s$ is subharmonic in $\mathcal{D}$.
\end{customthm}

The argument we gave at the end of the last proof applies generally to subharmonic functions defined on an arbitrary bounded domain. More precisely, we have the following theorem (that justifies the use of the word "subharmonic"):

\begin{theorem}
\label{thm:c2-4-5}
Suppose $s$ is a continuous function on the closure of a bounded domain $\mathcal{D}$. If $s$ is subharmonic in $\mathcal{D}$ and $u$ is a continuous function on $\bar{\mathcal{D}}$ that is harmonic in $\mathcal{D}$ and satisfies $s(x) \leq u(x)$ for all $x \in \mathcal{D} = \bar{\mathcal{D}} - \mathcal{D}$ then $s(x) \leq u(x)$ for all $x \in \bar{\mathcal{D}}$.
\end{theorem}

\textsc{Proof.}\ Consider the function $s - u$. If $s(x) - u(x) > 0$ for some $x \in \mathcal{D}$ let $c = \sup_{x \in \overline{\mathcal{D}}} \{s(x) - u(x)\} > 0$ and $F = \{x \in \overline{\mathcal{D}}; s(x) - u(x) = c\}$ ($c < \infty$ since $\overline{\mathcal{D}}$ is compact). Then $F$ is obviously closed. On the other hand, $-u$, being a harmonic function, satisfies the mean-value property; therefore,

\[
c=s(x)-u(x)\leq\frac{1}{\omega_{n-1}}\int\limits_{\Sigma}\left\{s(x+\rho t^{\prime})-u(x+\rho t^{\prime})\right\}d t^{\prime}
\]

whenever $x \in F$ and the sphere $\{t; |t - x| \leq \rho\}$ is contained in $\mathcal{D}$. As in the previous proof we can, therefore, conclude that $F$ is open. It follows that $F$ must be empty (since $\bar{\mathcal{D}}$ is connected and $s(x) - u(x) \leq 0 < c$ when $x \in \bar{\mathcal{D}} - \mathcal{D}$) and it follows that $s(x) - u(x) \leq 0$ for all $x \in \bar{\mathcal{D}}$.

Before proceeding further in our brief study of subharmonic functions it will be useful to present some examples:

(1) It is clear that a harmonic function is subharmonic. Moreover, it follows immediately from the definition that $|u|$ is subharmonic whenever u is harmonic.

(2) If $s$ is subharmonic and $\varphi$ is a nondecreasing convex function defined on an interval containing the range of $s$ then the composition $\varphi \circ s$ is subharmonic. This follows from Jensen's inequality\footnote{One of the many forms of Jensen's inequality is the following: Suppose $\mu$ is a finite measure on the space M, $f$ a $\mu$-integrable function and $\varphi$ a convex function whose domain includes the range of $f$, then 
\begin{equation}\tag{*}
\varphi\Biggl(\int\limits_{M}f(t)\;d\mu(t)/\mu(M)\Biggr)\;\leq\int\limits_{M}\varphi(f(t))\;d\mu(t)/\mu(M).
\end{equation}
 If $M$ is the interval $[x, y]$, $\mu(\{x\}) = 1 = \mu(\{y\})$, while $\mu(E) = 0$ whenever a Borel set $E$ contains neither $x$ nor $y$, and $f$ is the identity function $f(t) = t$, for $x \leq t \leq y$, then this inequality becomes the defining inequality \eqref{eq:c2-4-1} for convex functions. It is not hard, however, to obtain (*) from \eqref{eq:c2-4-1} by a simple limiting argument.} applied to the right member of inequality \eqref{eq:c2-4-2}:

\[
\varphi(s(x_{0}))\leq\varphi\biggl(\int_{\Sigma}s(x_{0}+r t^{\prime})d t^{\prime}/\omega_{n-1}\biggr)\leq\int_{\Sigma}\varphi(s(x_{0}+r t^{\prime}))d t^{\prime}/\omega_{n-1}.
\]

Since $\varphi(x) = x^p, p \geq 1$, defines a convex function for $x \geq 0$, this observation, combined with the first example, gives us the fact that $|u|^p$ is subharmonic whenever $u$ is subharmonic.

(3) It is not true, in general, that $|u|^p$ is subharmonic for $p < 1$ when $u$ is a harmonic function. For example, no power less than 1 of $|u|$ is subharmonic if $u$ is the harmonic function with values $u(x, y) = x$, for $(x, y) \in E_2$. In two dimensions, however, a considerably stronger result than that stated at the end of the previous example is true if we consider a harmonic function together with a harmonic conjugate\index{harmonic conjugate@harmonic conjugate}. \footnote{If $u$ is harmonic in a region $\mathcal{R} \subset E_2$ a harmonic function, $v$, such that $F = u + iv$ is analytic in $\mathcal{R}$ is called a harmonic conjugate of $u$. It is an elementary fact that, if $\mathcal{R}$ is simply connected, then such a harmonic conjugate exists and is unique up to an additive constant. In particular, a harmonic conjugate of $u$ exists in a neighborhood of each point of the region $\mathcal{R}$ where $u$ is harmonic.} More precisely, if $F = u + iv$ is an analytic function, then $|F|^p$ is subharmonic for all $p > 0$. We can show this in the following way: Let $\mathcal{R}$ be the set of points, $x + iy = z$, in the domain of $F$ such that $F(z) \neq 0$; then $s = \log |F|$ is harmonic in $\mathcal{R}$ and, a fortiori, subharmonic. Since $\varphi(x) = e^{px}$, $p > 0$, defines a convex function, the composition $\varphi \circ s = |F|^p$ is subharmonic in $\mathcal{R}$. The fact that $|F|^p$ is subharmonic in the domain of $F$ now follows from Theorem \ref{thm:c2-4-4}.

(4) If $s_1, s_2, \ldots, s_k$ are subharmonic in a domain $\mathcal{D}$ and $\alpha_1, \alpha_2, \ldots, \alpha_k$

are nonnegative real numbers, then the linear combination $\alpha_1 s_1 + \alpha_2 s_2 + \cdots + \alpha_k s_k$ is a subharmonic function in $\mathcal{D}$. Similarly $s(x) = \max_{x \in \mathcal{D}} \{s_1(x), s_2(x)\}$ defines a subharmonic function in $\mathcal{D}$ when $s_1$ and $s_2$ are subharmonic in this domain

If a harmonic function, $u$, majorizes the function $s$ in a domain we say that $u$ is a harmonic majorant\index{harmonic majorant@harmonic majorant} of $s$. If $u \leq h$ whenever $h$ is another harmonic majorant of $s$, we say that $u$ is the least harmonic majorant of $s$. Theorem \ref{thm:c2-4-5} tells us that in case $s$ is continuous on the closure $\mathcal{D}$, of the bounded domain $\mathcal{D}$ and is subharmonic in $\mathcal{D}$, then $u$ is a harmonic majorant of $s$ in $\mathcal{D}$ provided it has a continuous extension to $\mathcal{D}$ that majorizes $s$ on $\partial\mathcal{D}$. If $\mathcal{D}$ admits the solution of the Dirichlet problem, it follows immediately from this that the least harmonic majorant of $s$ is the continuous function on $\mathcal{D}$ that is harmonic in $\mathcal{D}$ and equals $s$ on $\partial\mathcal{D} = \mathcal{D} - \mathcal{D}$. We end this section by proving a related result for certain subharmonic functions defined in upper half-spaces. This result will be used (see Chapter III and VI), in conjunction with Theorem \ref{thm:c2-3-19} or Theorem \ref{thm:c2-3-24}, to obtain the existence of boundary values for certain classes of harmonic functions.

\begin{theorem}
\label{thm:c2-4-6}
Suppose s is a nonnegative subharmonic function defined in $E_{n+1}^{+}$ satisfying
\begin{equation}\label{eq:c2-4-7}
\int\limits_{E_{n}}[s(x,y)]^{q} dx\leq c^{q}<\infty
\end{equation}

for all $y > 0$, where $1 \leq q < \infty$ and $c$ is independent of $y > 0$. Then $s$ has a least harmonic majorant in $E_{n+1}^{+}$. Moreover, in case $q > 1$ this majorant is the Poisson integral of a function $f \in L^{q}(E_{n})$, with $\|f\|_{q} \leq c$; if $q = 1$ it is the Poisson--Stieltjes integral of a finite Borel measure on $E_{n}$ having total measure not exceeding $c$.
\end{theorem}

\textsc{Proof.}\ We first observe that $s$ satisfies the hypotheses of Lemma \ref{lem:c2-2-6} except for harmonicity. The proof of that lemma, however, depends only on inequality \eqref{eq:c2-4-2} and not on the full force of the mean-value property of harmonic functions. Consequently, there exists a constant $A = A_{n,p} > 0$ such that

\begin{equation}\label{eq:c2-4-8}
\left\Vert s(\cdot,y)\right\Vert_{\infty}=\sup_{x\in E_{n}}\left\vert s(x,y)\right\vert\leq A c y^{-n/q}.
\end{equation}

We shall show that, for each proper sub-half space $E_{n+1,y_0}^+$ = $(x, y) \in E_{n+1}$; $y \geq y_0 > 0$, $s$ belongs to the class $C_0(E_{n+1,y_0}^+)$. That is,

\begin{equation}\label{eq:c2-4-9}
\lim\limits_{|x|^{2}+y^{2}\to\infty}s(x,y)=0
\end{equation}

provided $y \geq y_0$. First we note that it follows from \eqref{eq:c2-4-8} that $s(x, y)$ is small provided $y$ is large. Thus, it suffices to show that $s(x, y)$ tends to 0 as $|x| \to \infty$ when $y$ lies between $y_0$ and some fixed $y_1 > y_0$. Let us fix a positive $r < y_0$ and put $B = \{(x, y) \in E_{n+1}; y_0 - r < y < y_1 + r\} \subset (E_{n+1}^+)$. Then

\[
\int_{B}[s(x,y)]^{q} dx\ dy=\int_{y_{0}-\tau}^{y_{1}+\tau}\int_{E_{n}}[s(x,y)]^{q} dx\ dy\ \leq(y_{1}+2r-y_{0})c^{q}<\infty.
\]

Consequently, if $B_k = B \cap \{(x, y) \in E_{n+1}; |x| \geq k\}$, we have

\begin{equation}\label{eq:c2-4-10}
\lim_{k\to\infty}\int\limits_{B_{k}}[s(x,y)]^{q} dx dy=0.
\end{equation}

If $|x| \geq m$ and $y_0 \leq y \leq y_1$, we can find a (solid) sphere $S$ of radius $r$ about the point $(x, y)$ that is contained in the set $B_{m-r}$. Since $\varphi(t) = t^q$ defines a nondecreasing convex function, $\varphi \circ s = s^q$ is also subharmonic (see example (2) above) and, therefore,

\[
[s(x,y)]^{q}\leq\frac{1}{|S|}\intop_{S}[s(\xi,\eta)]^{q}\mathop{}{d}\xi\mathop{}{d}\eta\leq\frac{1}{|S|}\intop_{B_{m-r}}[s(\xi,\eta)]^{q}\mathop{}{d}\xi\mathop{}{d}\eta.\footnote{Inequality \eqref{eq:c2-4-2} asserts that the mean value of s over the surface of a sphere of radius r about $x_{0}$ dominates $s(x_{0})$. By the same argument that was used in the proof of Lemma \ref{lem:c2-2-6} we can show that $s(x_{0})$ is dominated by the mean value of s over the solid sphere of radius r about $x_{0}$. It is this inequality, applied to $s^{q}$ at the point $(x, y)$, that is being used here.}
\]

But, by \eqref{eq:c2-4-10}, the last term tends to 0 as k = m - r tends to $\infty$, and \eqref{eq:c2-4-9} is established.

In order to construct the desired harmonic majorant we define, for each $\epsilon > 0$ and y > 0,

\[
m_{\varepsilon}(x,y)=\int\limits_{E_{n}}s(x-t,\varepsilon)P(t,y)d t=\int\limits_{E_{n}}s(t,\varepsilon)P(x-t,y)d t.
\]

Then, by inequality \eqref{eq:c2-2-2}, we have

\begin{equation}\label{eq:c2-4-11}
\int_{E_{n}}[m_{\varepsilon}(x,y)]^{q}dx\leq c^{q}
\end{equation}

for all $\varepsilon > 0$ and y > 0. Moreover, it follows from \eqref{eq:c2-4-9} and part (b) of Theorem \ref{thm:c2-2-1} that

\[
0=\lim_{y\to0}\|s(\cdot,\varepsilon)-m_{\varepsilon}(\cdot,y)\|_{\infty}=\lim_{y\to0}\left\{\sup_{x\in E_{n}}|s(x,\varepsilon)-m_{\varepsilon}(x,y)|\right\}.
\]

From this we can conclude that, if $\delta$ is a positive number, there exists a sufficiently small $y_0 > 0$ such that $s(x, \varepsilon + y_0) - m_\varepsilon(x, y_0) < \delta$ for all $x \in E_n$. By \eqref{eq:c2-4-9} we see that $s(x, \varepsilon + y) - m_\varepsilon(x, y) < \delta$ if either $|x|$ or $y$ is large, say $|x| \geq k$ or $y \geq y_1$. Thus, the inequality

\[
s(x,\varepsilon+y)-m_{\varepsilon}(x,y)<\delta
\]

holds on the boundary of the region $R = \{(x, y) \in E_{n+1}; |x| \leq k, y_0 \leq y \leq y_1\}$. An examination of the proof of the maximum property for harmonic functions shows that only the mean value inequality \eqref{eq:c2-4-2} was used; therefore, this property is satisfied by subharmonic functions. Applying it to the subharmonic function $s(x, \varepsilon + y) - m_\varepsilon(x, y) (-m)$, being harmonic, is subharmonic) and the region $R$ we see that the last inequality is valid throughout $R$. Letting $\delta \to 0$ and, accordingly, letting $k, y_1 \to \infty$ and $y_0 \to 0$ we obtain

\begin{equation}\label{eq:c2-4-12}
s(x,\varepsilon+y)\leq m_{\varepsilon}(x,y)
\end{equation}

for all $(x, y) \in E_{n+1}^+$.

Employing the same argument that was used in the proof of Theorem \ref{thm:c2-2-5} (with $u$ replaced by $s$ and $p$ replaced by $q$) we can find a null sequence $\{\epsilon_k\}$ and an $L^q$ function $f$, when $q > 1$, such that

\[
\begin{aligned}\lim_{k\to\infty}m_{\varepsilon_{k}}(x,y)&=\lim_{k\to\infty}\int\limits_{E_{n}}s(t,\varepsilon_{k})P(x-t,y)dt\\&=\int\limits_{E_{n}}f(t)P(x-t,y)dt=m(x,y)\end{aligned}
\]

while, when q = 1, instead of f there exists a finite Borel measure $\mu$, such that

\[
\begin{aligned}\lim_{k\to\infty}m_{\varepsilon_{k}}(x,y)&=\lim_{k\to\infty}\int\limits_{E_{n}}s(t,\varepsilon_{k})P(x-t,y)dt\\&=\int\limits_{E_{n}}P(x-t,y)d\mu(t)=m(x,y).\end{aligned}
\]

In either case, letting $\varepsilon = \varepsilon_{k}$ tend to zero in inequality \eqref{eq:c2-4-12}, we obtain

\[
s(x,y)\leq m(x,y).
\]

The function $m$ is the desired Poisson (or Poisson-Stieltjes) integral. The fact that it is the least harmonic majorant of $s$ follows from the obvious fact that $m_{\epsilon}(x,y)$ is the least harmonic majorant of $s(x,\varepsilon+y)$: for any harmonic majorant, $h$, of $s$ satisfies

\[
m(x,y)=\lim_{k\to\infty}m_{\varepsilon_{k}}(x,y)\leq\lim_{k\to\infty}h(x,\varepsilon_{k}+y)=h(x,y).
\]

\section{Further Results}

\frp{5}{1} At the beginning of this chapter we introduced the mean value $\mathcal{M}_{x,u}(r) = \mathcal{M}(r)$ of a function $u$ taken over the surface of the sphere $S_r(x)$ of radius $r$ centered at $x$. It is just as natural to form the mean value

\[
\mathcal{A}_{x,u}(r)=\mathcal{A}(r)=\frac{1}{\Omega_{n}r^{n}}\int_{|t|\leq r}u(x+t)\,dt,
\]

of the function $u$, taken over the volume of the sphere $S_r(x)$. We can, therefore, consider two different mean-value properties: the one introduced at the beginning of this chapter and the analogous one involving $\mathcal{A}(r)$. The argument used at the beginning of the proof of Lemma \ref{lem:c2-2-6} shows that if a function satisfies the first of these mean value properties then it satisfies the second. Similarly, as we have seen in the proof of Theorem \ref{thm:c2-4-6}, a function satisfying inequality \eqref{eq:c2-4-2} satisfies the inequality obtained by replacing $\mathcal{M}(r)$ by $\mathcal{A}(r)$. Moreover, it is not hard to show that, by replacing $\mathcal{M}(r)$ by $\mathcal{A}(r)$, we obtain another characterization of harmonic and subharmonic functions, similar to the ones in Theorem \ref{thm:c2-1-7} and inequality \eqref{eq:c2-4-2} (see Rado [1]).

\frp{5}{2} There are several results concerning harmonic (and analytic) functions defined on unbounded domains that are extensions of the maximum principle as stated in Corollary \ref{cor:c2-1-3p}. These results can be described in a general way as follows: Suppose we are given a continuous function, defined on a closed unbounded region $\mathcal{E} \subset E_n$, that is harmonic in the interior of $\mathcal{E}$. If the function satisfies certain weak restrictions on its growth (as $x \in \mathcal{E}$ tends to $\infty$) and satisfies certain more stringent conditions on the boundary of $\mathcal{E}$ (often the function is assumed to be bounded on $\partial \mathcal{E} = \bar{\mathcal{E}} - \mathcal{E}$), then it satisfies these more stringent conditions throughout $\mathcal{E}$. Such results are referred to as Phragmén\index{Phragmen-Lindelof theorems@Phragmen-Lindelof theorems}--Lindelöf theorems. Corollary \ref{cor:c2-1-15} is a particularly simple example of such a result. The following is an example of a Phragmén--Lindelöf theorem that is considerably more general: Suppose u is harmonic in $E_{n+1}$, continuous on $\overline{E_{n+1}} = E_{n+1} \cup E_n$ and satisfies $u(x, y) = o(e^{|x|})$ as $|x| \to \infty$, for all a > 0, in each strip of the form $S = \{(x, y) \in E_{n+1}; 0 < y \leq y_0\}$. If, in addition, $u(x, y) = o(y)$ as

(a) $u(x, 0) \leq A < \infty$ for all $x \in E_n$ implies $u(x, y) \leq A$ for all $(x, y) \in E_{n+1}^+$;

(b) $u(x,0) \geq B > -\infty$ for all $x \in E_n$ implies $u(x,y) \geq B$ for all $(x,y) \in E_{n+1}^+$;

(c) $|u(x,0)| \leq C < \infty$ for all $x \in E_n$ implies $|u(x,y)| \leq C$ for all $(x,y) \in E_{n+1}^+$.

To show, say, part (b) we can assume $B = 0$ since the general case can be reduced to this one by considering the function $u(x, y) - B$. Fix $\delta > 0$ and let $g(x, y) = \sin(a\sqrt{n}(y + \delta)) \prod_{k=1}^{n} \cosh ax_k$; then $g$ is a harmonic function in $E_{n+1}$ (see example (6) at the beginning of this chapter). Part (b) would be established if we prove that $u(x, y) + ey + \eta g(x, y) \geq 0$ in $\overline{E_{n+1}^+}$ for each fixed $\varepsilon > 0$ and $\eta > 0$. By the minimum principle we see that it suffices to show that this inequality holds on the boundary of a cylinder of the form $\{(x, y) \in E_{n+1}; |x| \leq N, 0 \leq y \leq y_0\}$ for all $N$ and $y_0$ that are sufficiently large. Choosing $a > 0$ and $\delta > 0$ so that $a < \pi/\sqrt{n}(y_0 + \delta)$, it follows that $\sin(a\sqrt{n}(y + \delta))$ is positive and bounded away from $0$ when $0 \leq y \leq y_0$ and, moreover,

\[
\prod_{k=1}^{n}\text{cosh}a x_{k}\geq2^{-n}\prod_{k=1}^{n}e^{a|x_{k}|}=2^{-n}\exp\bigg\{a\sum_{1}^{n}|x_{k}|\bigg\}\geq2^{-n}\exp\{a|x|\}.
\]

These estimates, together with the hypotheses $u(x, y) = o(e^{a|x|})$ and $u(x, y) = o(y)$, imply the desired inequality $u(x, y) + \varepsilon y + \eta g(x, y) \geq 0$ for $|x| \geq N$ (if N is sufficiently large) and $0 \leq y \leq y_0$. By decreasing a appropriately we can choose the number $y_0$ to be as large as we wish. This, then, implies (b); part (a) is obtained from (b) by considering -u, and these two results clearly imply (c).

\frp{5}{3} We showed (see Lemma \ref{lem:c2-2-6} and inequality \eqref{eq:c2-4-8}) that if a (non-negative) subharmonic function $s$ in $E_{n+1}^+$ satisfied $\|s(\cdot, y)\|_p \leq c$ for all $y > 0$ and some $p \geq 1$, then $\|s(\cdot, y)\|_{\infty} = O(y^{-n/p})$. As a consequence of this result we find, more generally, that $\|s(\cdot, y)\|_q = O(y^{(n/q) - (n/p)})$ when $q \geq p$. To see this we observe that,

\[
\begin{aligned}\int\limits_{E_{n}}[s(x,y)]^{q} dx&=\int\limits_{E_{n}}[s(x,y)]^{p}[s(x,y)]^{q-p}dx\leq\|s(\cdot,y)\|_{\infty}^{q-p}c^{p}\\&\leq A[y^{-n/p}]^{(q-p)},\end{aligned}
\]

and the desired inequality follows by taking q-th roots.

\frp{5}{4} Theorems \ref{thm:c2-3-19} and 3.24 indicate that the concept of nontangential convergence at the boundary is essential. That this is indeed the fact can be seen from the following two results dealing with functions defined on the unit disk (Analogous results for the upper half-plane can be obtained with the aid of conformal mappings):

(i) Let $C_0$ be a simple closed curve passing through the point $z = 1$, all of whose other points are contained in the interior of the unit circle $K = \{z = x + iy; |z| = 1\}$. Suppose, further, that $C_0$ is tangent to the unit circle at 1 and that $C_\theta$ is the curve obtained by rotating $C_0$ about the origin through the angle $\theta$. Then there exists a bounded analytic function,

$F$, defined in the interior of $K$ such that, for almost all $\theta$, $F(z)$ does not tend to any limit as $z$ tends to $e^{i\theta}$ within $C_\theta$ (see Zygmund [1], Vol. I, page 280).

(ii) Let G be a continuous function defined in the interior of the unit circle, K, and let E be a set of the first category in K. Then there exists an analytic function F defined in the interior of K such that

\[
\lim_{r\to1}\left\{F(r e^{i\theta})\:-\:G(r e^{i\theta})\right\}=0
\]

whenever $e^{i\theta} \in E$ (see Zygmund [1], Vol. II, page 204).

Since $E$ can have measure $2\pi$, this last result asserts that the radial behavior, at almost all points of $K$, of analytic functions is no better than that of continuous functions. Thus, for example, an analytic function $F$ can tend to 0 along almost every radius and, yet, not vanish identically. On the other hand, an analytic function having nontangential limit 0 at the points of a set of positive measure in $K$ must be identically zero (see Zygmund [1], Vol. II, page 203). Moreover, an analytic function can be bounded (though, obviously, not uniformly) on almost every radius and, yet, fail to have radial limits at almost all points of $K$.

\frp{5}{5} Theorem \ref{thm:c2-3-19} has several extensions. Carleson (see [1]) has shown that the same result holds if we assume that the (real-valued) harmonic function $u$ is nontangentially bounded from below or from above at each point of $S$. In the two-dimensional case this result can easily be extended to regions more general than the upper half-plane with the aid of conformal mappings. Hunt and Wheeden (see [1]) have obtained a similar extension in higher dimensions for a domain $\mathcal{D}$ whose boundary satisfies a Lipschitz condition. They have shown that if $u$ is harmonic in such a domain $\mathcal{D}$, $S \subset \partial \mathcal{D}$ has the property that for each $q \in S$ there exists an open cone with vertex $q$ that is exterior to $\mathcal{D}$, and $u$ is bounded from below (or above) in a truncated cone having vertex $q$, for each $q \in S$, then $u$ has a finite nontangential limit at each point of $S$ except, possibly, for a set of harmonic measure zero.

\frp{5}{6} By slightly modifying the argument used in the proof of Lemma \ref{lem:c2-3-3} we can obtain the following closely related result: Suppose that to each point $x$ of the bounded set $F \subset E_n$ there corresponds a non-empty open sphere $S_x$ with center $x$. Then there exist a countable collection $\{S_{x_i}\}$ of these spheres covering $F$ and such that each $x \in E_n$ belongs to at most $a_n$ of these spheres ( $a_n$ depends only on the dimension; in fact, it can be shown that $a_n$ can be chosen to be $3^n$). Using this result we can generalize Theorem \ref{thm:c2-3-4}: Suppose $\nu_1$ and $\nu_2$ are two nonnegative finite Borel measures on $E_n$ and $\varphi(x) = \sup\{\nu_1(S_x)/\nu_2(S_x)\}$, where the supremum is taken over all open spheres $S_x$ centered at $x \in E_n$ (when $\nu_2$ is Lebesgue

measure and $\nu_1(S_x) = \int_{S_x} f$, for some integrable $f$, then $\varphi$ is the Hardy--Littlewood maximal function $m_f$. Then, if $F_s = \{x \in E_n: \varphi(x) > s > 0\}$, we have $\nu_2(F_s) \leq a_n \nu_1(E_n) s^{-1}$. From this one can show that $\nu_1$ is differentiable with respect to $\nu_2$; that is, $\lim \{\nu_1(S_x)/\nu_2(S_x)\} = f(x)$ exists for all $x \in E_n$ outside a set of $\nu_2$-measure zero, the limit being taken over spheres centered at $x$ with radii tending to zero. If we decompose $\nu_1$ into its absolutely continuous and singular parts, with respect to $\nu_2$, then the Radon--Nikodym derivative of the former equals $f(x)$ almost everywhere (again, with respect to $\nu_2$). Finally, we remark that these results still hold if other geometric figures are allowed to replace spheres (see Cotlar [1], pp. 126--127).

\frp{5}{7} In theorem 3.22 and Corollary \ref{cor:c2-3-23} we assumed $1 < p < \infty$. The case $p = \infty$ was discussed at the end of the proof of Theorem \ref{thm:c2-3-22}. There, we pointed out that the result in this case follows from similar arguments. The case $p = 1$ is quite different: one can show that there exist functions $f \in L^1(E_2)$, e.g., such that the iterated Poisson integrals

\[
u(x_{1},\varepsilon_{1},x_{2},\varepsilon_{2})=\int\limits_{-\infty}^{\infty}\int\limits_{-\infty}^{\infty}f(x_{1}-t_{1},x_{2}-t_{2})P(t_{1},\varepsilon_{1})P(t_{2},\varepsilon_{2})dt_{1}dt_{2}
\]

do not converge to $f(x_1, x_2)$ as $\epsilon_1, \epsilon_2 \to 0$ for almost all $x = (x_1, x_2) \in E_2$ (see Jessen, Marcinkiewicz, and Zygmund [1] where the iterated Poisson integral for the unit circle is discussed. The argument used there can be modified to include the case under consideration here). See also \eqref{frp:c3-6-6} and \eqref{frp:c3-6-8} of Chapter III.

\frp{5}{8} It follows easily from our development of subharmonic functions that if $s$ has continuous second partial derivatives then the following three conditions are equivalent: (i) $s$ is subharmonic; (ii) $\Delta s \geq 0$; (iii) if $\mathcal{R}$ is a bounded domain whose closure is contained in the domain of definition of $s$, $u$ is a continuous function on $\tilde{\mathcal{R}}$ that is harmonic in $\mathcal{R}$ and satisfies $u \geq s$ on the boundary $\partial \mathcal{R} = \tilde{\mathcal{R}} - \mathcal{R}$, then $u \geq s$ on $\mathcal{R}$ (see Theorem \ref{thm:c2-4-5}). There are several ways of extending these results. We shall sketch one such extension that is closely connected with the ideas developed in our treatment of tempered distributions (see Section 3 of Chapter I). Let us say that a continuous function $s$, defined on some domain $D$, had a nonnegative generalized Laplacian if $\int_{D} s \Delta \varphi \geq 0$ whenever $\varphi$ is a nonnegative $C^{\infty}$ function with compact support interior to $D$ (integrating by parts, we see that if $s$ has continuous second partial derivatives then this property is equivalent to (ii) above). We first claim that such a function $s$ is subharmonic. In order to do this we show that $s$ satisfies the mean-value inequality \eqref{eq:c2-4-2}. In view of Theorem \ref{thm:c1-1-25}, Chapter

I, it suffices to show that this is the case for each of the (everywhere defined) functions $s_\varepsilon(x) = \int_D s(t)\varphi_\varepsilon(x - t)\,dt$, where $\varphi$ is a nonnegative $C^\infty$ function with compact support having integral 1 and $\varepsilon > 0$ sufficiently small. By Theorem \ref{thm:c2-4-3} we see that this is true provided $\Delta s_\varepsilon \geq 0$. As we have observed above, this is equivalent to $\int_D s_\varepsilon\Delta\psi \geq 0$ for all nonnegative $C^\infty$ functions $\psi$ having compact support within $D$. Choosing $\varepsilon$ small and the support of $\psi$ within an appropriate compact subset of $D$ we have $\int_D s_\varepsilon\Delta\psi = \int_D s\Delta(\tilde{\varphi}_\varepsilon \cdot \psi) \geq 0$ ($\tilde{\varphi}$ denotes the reflection of $\varphi$, as defined in Chapter I preceding equality \eqref{eq:c1-3-12}). This shows that $s$ is subharmonic. Theorem \ref{thm:c2-4-5} tells us that (i), in turn, implies the harmonic majorization property (iii). It is immediate, however, that the latter implies the mean-value inequality \eqref{eq:c2-4-2}. Assuming that $s$ satisfies this last property we shall show that it has a nonnegative generalized Laplacian. Let $\eta_\delta$ be the function obtained after dividing by $\Omega_n\delta^n$ the characteristic function of the sphere of radius $\delta$ about 0. Then \eqref{eq:c2-4-2} can be expressed as the inequality $s \leq s \star \eta_\delta$. Thus, $\int_D s\varphi \leq \int_D (s \star \eta_\delta)\varphi$ or, equivalently,

\[
\int\limits_{D}s(\eta_{\delta}*\varphi-\varphi)\geq0
\]

for all nonnegative $C^\infty$ functions with compact support interior to $D$. But $\{\eta_\delta*\varphi-\varphi\}/\delta^2$ tends to a positive constant times $\Delta\varphi$ as $\delta$ tends to $0$; therefore, $\int_D s\,\Delta\varphi \geq 0$.

\bibnoteshead

Kellog [1] is a comprehensive text on the theory of harmonic functions of three variables. For a more modern treatment of the general $n$-dimensional case see Brelot [1]. The characterization of Poisson integrals is a classical result and the proofs used here are not very different from those used for the characterization of the Cesáro and Abel means of trigonometric series that can be found in Zygmund [1], Chapter III. Theorem \ref{thm:c2-3-10} and inequality \eqref{eq:c2-3-18} are an extension to upper half-spaces of the corresponding results for Poisson integrals of functions defined on the unit circle originally obtained by Hardy and Littlewood [1]; for these and other variants see K. T. Smith [1]. Theorems \ref{thm:c2-3-19} and 3.24 are due to Calderón [1]. The maximal function was first introduced, in one dimension, by Hardy and Littlewood [1]. The extension to several variables is due to Wiener [2]. In connection with the covering lemma used here see Wiener [2], Marcinkiewicz and Zygmund [1], Besicovitch [2], Cotlar [1], and de Guzmán [1]. Theorem \ref{thm:c2-3-12} is closely related to a result of Cotlar [1], an earlier and more general result is due to Banach [1]. For a comprehensive treatment of the theory of subharmonic functions of two variables see Rado [1]. Theorem \ref{thm:c2-4-6} can be found in Stein and G. Weiss [1]. For a more complete treatment of strong differentiability of integrals see Saks [1].

